%% Ngân hàng đề Toán 12 (KNTT) -- Chương IV, Bài 13, Dạng 02
%% Dạng: Diện tích hình phẳng giới hạn bởi hai đồ thị
%% Phương án ĐẦU TIÊN trong \cauTN luôn là đáp án đúng (gói trondethi tự đảo).
%% Phần II: \y{nội dung}{d|s}. Phần III: đáp án đặt ở đối số cuối.

\cauTN[id=C4-B13-D02-NB-001,muc=NB,loigiai={Theo công thức: $S=\displaystyle\int_{1}^{3}\left|f(x)-g(x)\right|\,\mathrm{d}x$.}]{Cho các hàm số $y=f(x)$, $y=g(x)$ liên tục trên đoạn $[1;3]$. Diện tích $S$ của hình phẳng giới hạn bởi đồ thị hai hàm số đó và hai đường thẳng $x=1$, $x=3$ được tính theo công thức}
  {$S=\displaystyle\int_{1}^{3}\left|f(x)-g(x)\right|\,\mathrm{d}x$}
  {$S=\pi\displaystyle\int_{1}^{3}\left|f(x)-g(x)\right|\,\mathrm{d}x$}{$S=\displaystyle\int_{1}^{3}\left[f(x)-g(x)\right]\,\mathrm{d}x$}{$S=\pi\displaystyle\int_{1}^{3}\left[f^2(x)-g^2(x)\right]\,\mathrm{d}x$}

\cauTN[id=C4-B13-D02-NB-002,muc=NB,loigiai={Phương trình hoành độ giao điểm: $x^{2} - 2 x + 1=x - 1\Leftrightarrow x^{2} - 3 x + 2=0\Leftrightarrow x\in\{1;2\}$. $S=\displaystyle\int_{1}^{2}\left|x^{2} - 3 x + 2\right|\,\mathrm{d}x=\dfrac{1}{6}$.}]{Diện tích hình phẳng giới hạn bởi parabol $y=x^{2} - 2 x + 1$ và đường thẳng $y=x - 1$ bằng}
  {$\dfrac{1}{6}$}
  {$-\dfrac{1}{6}$}{$\dfrac{\pi}{6}$}{$\dfrac{1}{3}$}

\cauTN[id=C4-B13-D02-NB-003,muc=NB,loigiai={Phương trình hoành độ giao điểm: $x^{2} + 3 x - 5=2 x + 1\Leftrightarrow x^{2} + x - 6=0\Leftrightarrow x\in\{-3;2\}$. $S=\displaystyle\int_{-3}^{2}\left|x^{2} + x - 6\right|\,\mathrm{d}x=\dfrac{125}{6}$.}]{Diện tích hình phẳng giới hạn bởi parabol $y=x^{2} + 3 x - 5$ và đường thẳng $y=2 x + 1$ bằng}
  {$\dfrac{125}{6}$}
  {$-\dfrac{125}{6}$}{$\dfrac{125 \pi}{6}$}{$\dfrac{125}{3}$}

\cauTN[id=C4-B13-D02-NB-004,muc=NB,loigiai={Phương trình hoành độ giao điểm: $- x^{2} - 6 x - 8=- x - 2\Leftrightarrow - x^{2} - 5 x - 6=0\Leftrightarrow x\in\{-3;-2\}$. $S=\displaystyle\int_{-3}^{-2}\left|- x^{2} - 5 x - 6\right|\,\mathrm{d}x=\dfrac{1}{6}$.}]{Diện tích hình phẳng giới hạn bởi parabol $y=- x^{2} - 6 x - 8$ và đường thẳng $y=- x - 2$ bằng}
  {$\dfrac{1}{6}$}
  {$-\dfrac{1}{6}$}{$\dfrac{\pi}{6}$}{$\dfrac{1}{3}$}

\cauTN[id=C4-B13-D02-NB-005,muc=NB,loigiai={Phương trình hoành độ giao điểm: $x^{2} + 3=3 - x\Leftrightarrow x^{2} + x=0\Leftrightarrow x\in\{-1;0\}$. $S=\displaystyle\int_{-1}^{0}\left|x^{2} + x\right|\,\mathrm{d}x=\dfrac{1}{6}$.}]{Diện tích hình phẳng giới hạn bởi parabol $y=x^{2} + 3$ và đường thẳng $y=3 - x$ bằng}
  {$\dfrac{1}{6}$}
  {$-\dfrac{1}{6}$}{$\dfrac{\pi}{6}$}{$\dfrac{1}{3}$}

\cauTN[id=C4-B13-D02-NB-006,muc=NB,loigiai={Phương trình hoành độ giao điểm: $x^{2} + 4 x + 5=x + 3\Leftrightarrow x^{2} + 3 x + 2=0\Leftrightarrow x\in\{-2;-1\}$. $S=\displaystyle\int_{-2}^{-1}\left|x^{2} + 3 x + 2\right|\,\mathrm{d}x=\dfrac{1}{6}$.}]{Diện tích hình phẳng giới hạn bởi parabol $y=x^{2} + 4 x + 5$ và đường thẳng $y=x + 3$ bằng}
  {$\dfrac{1}{6}$}
  {$-\dfrac{1}{6}$}{$\dfrac{\pi}{6}$}{$\dfrac{1}{3}$}

\cauTN[id=C4-B13-D02-NB-007,muc=NB,loigiai={Phương trình hoành độ giao điểm: $- x^{2} - 2 x - 1=x + 1\Leftrightarrow - x^{2} - 3 x - 2=0\Leftrightarrow x\in\{-2;-1\}$. $S=\displaystyle\int_{-2}^{-1}\left|- x^{2} - 3 x - 2\right|\,\mathrm{d}x=\dfrac{1}{6}$.}]{Diện tích hình phẳng giới hạn bởi parabol $y=- x^{2} - 2 x - 1$ và đường thẳng $y=x + 1$ bằng}
  {$\dfrac{1}{6}$}
  {$-\dfrac{1}{6}$}{$\dfrac{\pi}{6}$}{$\dfrac{1}{3}$}

\cauTN[id=C4-B13-D02-NB-008,muc=NB,loigiai={Theo công thức: $S=\displaystyle\int_{a}^{b}\left|f(x)-g(x)\right|\,\mathrm{d}x$.}]{Cho các hàm số $y=f(x)$, $y=g(x)$ liên tục trên đoạn $[a;b]$. Diện tích $S$ của hình phẳng giới hạn bởi đồ thị hai hàm số đó và hai đường thẳng $x=a$, $x=b$ được tính theo công thức}
  {$S=\displaystyle\int_{a}^{b}\left|f(x)-g(x)\right|\,\mathrm{d}x$}
  {$S=\displaystyle\int_{a}^{b}\left|f(x)+g(x)\right|\,\mathrm{d}x$}{$S=\pi\displaystyle\int_{a}^{b}\left[f^2(x)-g^2(x)\right]\,\mathrm{d}x$}{$S=\displaystyle\int_{a}^{b}\left[f(x)-g(x)\right]\,\mathrm{d}x$}

\cauTN[id=C4-B13-D02-NB-009,muc=NB,loigiai={Phương trình hoành độ giao điểm: $x^{2} + 5 x=x - 3\Leftrightarrow x^{2} + 4 x + 3=0\Leftrightarrow x\in\{-3;-1\}$. $S=\displaystyle\int_{-3}^{-1}\left|x^{2} + 4 x + 3\right|\,\mathrm{d}x=\dfrac{4}{3}$.}]{Diện tích hình phẳng giới hạn bởi parabol $y=x^{2} + 5 x$ và đường thẳng $y=x - 3$ bằng}
  {$\dfrac{4}{3}$}
  {$-\dfrac{4}{3}$}{$\dfrac{4 \pi}{3}$}{$\dfrac{8}{3}$}

\cauTN[id=C4-B13-D02-NB-010,muc=NB,loigiai={Phương trình hoành độ giao điểm: $- x^{2} - x + 5=-1\Leftrightarrow - x^{2} - x + 6=0\Leftrightarrow x\in\{-3;2\}$. $S=\displaystyle\int_{-3}^{2}\left|- x^{2} - x + 6\right|\,\mathrm{d}x=\dfrac{125}{6}$.}]{Diện tích hình phẳng giới hạn bởi parabol $y=- x^{2} - x + 5$ và đường thẳng $y=-1$ bằng}
  {$\dfrac{125}{6}$}
  {$-\dfrac{125}{6}$}{$\dfrac{125 \pi}{6}$}{$\dfrac{125}{3}$}

\cauTN[id=C4-B13-D02-NB-011,muc=NB,loigiai={Phương trình hoành độ giao điểm: $- x^{2} - 1=- x - 3\Leftrightarrow - x^{2} + x + 2=0\Leftrightarrow x\in\{-1;2\}$. $S=\displaystyle\int_{-1}^{2}\left|- x^{2} + x + 2\right|\,\mathrm{d}x=\dfrac{9}{2}$.}]{Diện tích hình phẳng giới hạn bởi parabol $y=- x^{2} - 1$ và đường thẳng $y=- x - 3$ bằng}
  {$\dfrac{9}{2}$}
  {$-\dfrac{9}{2}$}{$\dfrac{9 \pi}{2}$}{$9$}

\cauTN[id=C4-B13-D02-NB-012,muc=NB,loigiai={Phương trình hoành độ giao điểm: $x^{2} - 4 x - 3=- x - 3\Leftrightarrow x^{2} - 3 x=0\Leftrightarrow x\in\{0;3\}$. $S=\displaystyle\int_{0}^{3}\left|x^{2} - 3 x\right|\,\mathrm{d}x=\dfrac{9}{2}$.}]{Diện tích hình phẳng giới hạn bởi parabol $y=x^{2} - 4 x - 3$ và đường thẳng $y=- x - 3$ bằng}
  {$\dfrac{9}{2}$}
  {$-\dfrac{9}{2}$}{$\dfrac{9 \pi}{2}$}{$9$}

\cauTN[id=C4-B13-D02-NB-013,muc=NB,loigiai={Theo công thức: $S=\displaystyle\int_{-1}^{1}\left|f_1(x)-f_2(x)\right|\,\mathrm{d}x$.}]{Cho các hàm số $y=f_1(x)$, $y=f_2(x)$ liên tục trên đoạn $[-1;1]$. Diện tích $S$ của hình phẳng giới hạn bởi đồ thị hai hàm số đó và hai đường thẳng $x=-1$, $x=1$ được tính theo công thức}
  {$S=\displaystyle\int_{-1}^{1}\left|f_1(x)-f_2(x)\right|\,\mathrm{d}x$}
  {$S=\displaystyle\int_{-1}^{1}\left|f_1(x)+f_2(x)\right|\,\mathrm{d}x$}{$S=\displaystyle\int_{-1}^{1}\left[f_1(x)-f_2(x)\right]\,\mathrm{d}x$}{$S=\pi\displaystyle\int_{-1}^{1}\left[f_1^2(x)-f_2^2(x)\right]\,\mathrm{d}x$}

\cauTN[id=C4-B13-D02-NB-014,muc=NB,loigiai={Phương trình hoành độ giao điểm: $x^{2} - 2 x + 2=2 x - 1\Leftrightarrow x^{2} - 4 x + 3=0\Leftrightarrow x\in\{1;3\}$. $S=\displaystyle\int_{1}^{3}\left|x^{2} - 4 x + 3\right|\,\mathrm{d}x=\dfrac{4}{3}$.}]{Diện tích hình phẳng giới hạn bởi parabol $y=x^{2} - 2 x + 2$ và đường thẳng $y=2 x - 1$ bằng}
  {$\dfrac{4}{3}$}
  {$-\dfrac{4}{3}$}{$\dfrac{4 \pi}{3}$}{$\dfrac{8}{3}$}

\cauTN[id=C4-B13-D02-NB-015,muc=NB,loigiai={Phương trình hoành độ giao điểm: $- x^{2} + 2 x + 7=2 x + 3\Leftrightarrow 4 - x^{2}=0\Leftrightarrow x\in\{-2;2\}$. $S=\displaystyle\int_{-2}^{2}\left|4 - x^{2}\right|\,\mathrm{d}x=\dfrac{32}{3}$.}]{Diện tích hình phẳng giới hạn bởi parabol $y=- x^{2} + 2 x + 7$ và đường thẳng $y=2 x + 3$ bằng}
  {$\dfrac{32}{3}$}
  {$-\dfrac{32}{3}$}{$\dfrac{32 \pi}{3}$}{$\dfrac{64}{3}$}

\cauTN[id=C4-B13-D02-NB-016,muc=NB,loigiai={Phương trình hoành độ giao điểm: $x^{2} - 3=- x - 1\Leftrightarrow x^{2} + x - 2=0\Leftrightarrow x\in\{-2;1\}$. $S=\displaystyle\int_{-2}^{1}\left|x^{2} + x - 2\right|\,\mathrm{d}x=\dfrac{9}{2}$.}]{Diện tích hình phẳng giới hạn bởi parabol $y=x^{2} - 3$ và đường thẳng $y=- x - 1$ bằng}
  {$\dfrac{9}{2}$}
  {$-\dfrac{9}{2}$}{$\dfrac{9 \pi}{2}$}{$9$}

\cauTN[id=C4-B13-D02-NB-017,muc=NB,loigiai={Theo công thức: $S=\displaystyle\int_{a}^{b}\left|f_1(x)-f_2(x)\right|\,\mathrm{d}x$.}]{Cho các hàm số $y=f_1(x)$, $y=f_2(x)$ liên tục trên đoạn $[a;b]$. Diện tích $S$ của hình phẳng giới hạn bởi đồ thị hai hàm số đó và hai đường thẳng $x=a$, $x=b$ được tính theo công thức}
  {$S=\displaystyle\int_{a}^{b}\left|f_1(x)-f_2(x)\right|\,\mathrm{d}x$}
  {$S=\pi\displaystyle\int_{a}^{b}\left[f_1^2(x)-f_2^2(x)\right]\,\mathrm{d}x$}{$S=\displaystyle\int_{a}^{b}\left[f_1(x)-f_2(x)\right]\,\mathrm{d}x$}{$S=\pi\displaystyle\int_{a}^{b}\left|f_1(x)-f_2(x)\right|\,\mathrm{d}x$}

\cauTN[id=C4-B13-D02-NB-018,muc=NB,loigiai={Phương trình hoành độ giao điểm: $- x^{2} + x + 1=- x - 2\Leftrightarrow - x^{2} + 2 x + 3=0\Leftrightarrow x\in\{-1;3\}$. $S=\displaystyle\int_{-1}^{3}\left|- x^{2} + 2 x + 3\right|\,\mathrm{d}x=\dfrac{32}{3}$.}]{Diện tích hình phẳng giới hạn bởi parabol $y=- x^{2} + x + 1$ và đường thẳng $y=- x - 2$ bằng}
  {$\dfrac{32}{3}$}
  {$-\dfrac{32}{3}$}{$\dfrac{32 \pi}{3}$}{$\dfrac{64}{3}$}

\cauTN[id=C4-B13-D02-NB-019,muc=NB,loigiai={Phương trình hoành độ giao điểm: $x^{2} - 3=- x - 3\Leftrightarrow x^{2} + x=0\Leftrightarrow x\in\{-1;0\}$. $S=\displaystyle\int_{-1}^{0}\left|x^{2} + x\right|\,\mathrm{d}x=\dfrac{1}{6}$.}]{Diện tích hình phẳng giới hạn bởi parabol $y=x^{2} - 3$ và đường thẳng $y=- x - 3$ bằng}
  {$\dfrac{1}{6}$}
  {$-\dfrac{1}{6}$}{$\dfrac{\pi}{6}$}{$\dfrac{1}{3}$}

\cauTN[id=C4-B13-D02-NB-020,muc=NB,loigiai={Theo công thức: $S=\displaystyle\int_{1}^{3}\left|f_1(x)-f_2(x)\right|\,\mathrm{d}x$.}]{Cho các hàm số $y=f_1(x)$, $y=f_2(x)$ liên tục trên đoạn $[1;3]$. Diện tích $S$ của hình phẳng giới hạn bởi đồ thị hai hàm số đó và hai đường thẳng $x=1$, $x=3$ được tính theo công thức}
  {$S=\displaystyle\int_{1}^{3}\left|f_1(x)-f_2(x)\right|\,\mathrm{d}x$}
  {$S=\displaystyle\int_{1}^{3}\left[f_1(x)-f_2(x)\right]\,\mathrm{d}x$}{$S=\pi\displaystyle\int_{1}^{3}\left|f_1(x)-f_2(x)\right|\,\mathrm{d}x$}{$S=\displaystyle\int_{1}^{3}\left|f_1(x)+f_2(x)\right|\,\mathrm{d}x$}

\cauTN[id=C4-B13-D02-NB-021,muc=NB,loigiai={Phương trình hoành độ giao điểm: $3 - x^{2}=x + 1\Leftrightarrow - x^{2} - x + 2=0\Leftrightarrow x\in\{-2;1\}$. $S=\displaystyle\int_{-2}^{1}\left|- x^{2} - x + 2\right|\,\mathrm{d}x=\dfrac{9}{2}$.}]{Diện tích hình phẳng giới hạn bởi parabol $y=3 - x^{2}$ và đường thẳng $y=x + 1$ bằng}
  {$\dfrac{9}{2}$}
  {$-\dfrac{9}{2}$}{$\dfrac{9 \pi}{2}$}{$9$}

\cauTN[id=C4-B13-D02-NB-022,muc=NB,loigiai={Phương trình hoành độ giao điểm: $x^{2} + 7 x + 8=2 x + 2\Leftrightarrow x^{2} + 5 x + 6=0\Leftrightarrow x\in\{-3;-2\}$. $S=\displaystyle\int_{-3}^{-2}\left|x^{2} + 5 x + 6\right|\,\mathrm{d}x=\dfrac{1}{6}$.}]{Diện tích hình phẳng giới hạn bởi parabol $y=x^{2} + 7 x + 8$ và đường thẳng $y=2 x + 2$ bằng}
  {$\dfrac{1}{6}$}
  {$-\dfrac{1}{6}$}{$\dfrac{\pi}{6}$}{$\dfrac{1}{3}$}

\cauTN[id=C4-B13-D02-NB-023,muc=NB,loigiai={Theo công thức: $S=\displaystyle\int_{1}^{3}\left|u(x)-v(x)\right|\,\mathrm{d}x$.}]{Cho các hàm số $y=u(x)$, $y=v(x)$ liên tục trên đoạn $[1;3]$. Diện tích $S$ của hình phẳng giới hạn bởi đồ thị hai hàm số đó và hai đường thẳng $x=1$, $x=3$ được tính theo công thức}
  {$S=\displaystyle\int_{1}^{3}\left|u(x)-v(x)\right|\,\mathrm{d}x$}
  {$S=\displaystyle\int_{1}^{3}\left|u(x)+v(x)\right|\,\mathrm{d}x$}{$S=\pi\displaystyle\int_{1}^{3}\left|u(x)-v(x)\right|\,\mathrm{d}x$}{$S=\displaystyle\int_{1}^{3}\left[u(x)-v(x)\right]\,\mathrm{d}x$}

\cauTN[id=C4-B13-D02-NB-024,muc=NB,loigiai={Phương trình hoành độ giao điểm: $- x^{2} + 5 x - 3=3\Leftrightarrow - x^{2} + 5 x - 6=0\Leftrightarrow x\in\{2;3\}$. $S=\displaystyle\int_{2}^{3}\left|- x^{2} + 5 x - 6\right|\,\mathrm{d}x=\dfrac{1}{6}$.}]{Diện tích hình phẳng giới hạn bởi parabol $y=- x^{2} + 5 x - 3$ và đường thẳng $y=3$ bằng}
  {$\dfrac{1}{6}$}
  {$-\dfrac{1}{6}$}{$\dfrac{\pi}{6}$}{$\dfrac{1}{3}$}

\cauTN[id=C4-B13-D02-NB-025,muc=NB,loigiai={Phương trình hoành độ giao điểm: $- x^{2} - x + 1=x - 2\Leftrightarrow - x^{2} - 2 x + 3=0\Leftrightarrow x\in\{-3;1\}$. $S=\displaystyle\int_{-3}^{1}\left|- x^{2} - 2 x + 3\right|\,\mathrm{d}x=\dfrac{32}{3}$.}]{Diện tích hình phẳng giới hạn bởi parabol $y=- x^{2} - x + 1$ và đường thẳng $y=x - 2$ bằng}
  {$\dfrac{32}{3}$}
  {$-\dfrac{32}{3}$}{$\dfrac{32 \pi}{3}$}{$\dfrac{64}{3}$}

\cauTN[id=C4-B13-D02-NB-026,muc=NB,loigiai={Theo công thức: $S=\displaystyle\int_{0}^{2}\left|u(x)-v(x)\right|\,\mathrm{d}x$.}]{Cho các hàm số $y=u(x)$, $y=v(x)$ liên tục trên đoạn $[0;2]$. Diện tích $S$ của hình phẳng giới hạn bởi đồ thị hai hàm số đó và hai đường thẳng $x=0$, $x=2$ được tính theo công thức}
  {$S=\displaystyle\int_{0}^{2}\left|u(x)-v(x)\right|\,\mathrm{d}x$}
  {$S=\displaystyle\int_{0}^{2}\left|u(x)+v(x)\right|\,\mathrm{d}x$}{$S=\pi\displaystyle\int_{0}^{2}\left|u(x)-v(x)\right|\,\mathrm{d}x$}{$S=\pi\displaystyle\int_{0}^{2}\left[u^2(x)-v^2(x)\right]\,\mathrm{d}x$}

\cauTN[id=C4-B13-D02-NB-027,muc=NB,loigiai={Phương trình hoành độ giao điểm: $x^{2} - 2 x - 6=3 - 2 x\Leftrightarrow x^{2} - 9=0\Leftrightarrow x\in\{-3;3\}$. $S=\displaystyle\int_{-3}^{3}\left|x^{2} - 9\right|\,\mathrm{d}x=36$.}]{Diện tích hình phẳng giới hạn bởi parabol $y=x^{2} - 2 x - 6$ và đường thẳng $y=3 - 2 x$ bằng}
  {$36$}
  {$-36$}{$36 \pi$}{$72$}

\cauTN[id=C4-B13-D02-NB-028,muc=NB,loigiai={Theo công thức: $S=\displaystyle\int_{0}^{2}\left|f(x)-g(x)\right|\,\mathrm{d}x$.}]{Cho các hàm số $y=f(x)$, $y=g(x)$ liên tục trên đoạn $[0;2]$. Diện tích $S$ của hình phẳng giới hạn bởi đồ thị hai hàm số đó và hai đường thẳng $x=0$, $x=2$ được tính theo công thức}
  {$S=\displaystyle\int_{0}^{2}\left|f(x)-g(x)\right|\,\mathrm{d}x$}
  {$S=\displaystyle\int_{0}^{2}\left|f(x)+g(x)\right|\,\mathrm{d}x$}{$S=\pi\displaystyle\int_{0}^{2}\left[f^2(x)-g^2(x)\right]\,\mathrm{d}x$}{$S=\pi\displaystyle\int_{0}^{2}\left|f(x)-g(x)\right|\,\mathrm{d}x$}

\cauTN[id=C4-B13-D02-NB-029,muc=NB,loigiai={Phương trình hoành độ giao điểm: $- x^{2} - 1=x - 3\Leftrightarrow - x^{2} - x + 2=0\Leftrightarrow x\in\{-2;1\}$. $S=\displaystyle\int_{-2}^{1}\left|- x^{2} - x + 2\right|\,\mathrm{d}x=\dfrac{9}{2}$.}]{Diện tích hình phẳng giới hạn bởi parabol $y=- x^{2} - 1$ và đường thẳng $y=x - 3$ bằng}
  {$\dfrac{9}{2}$}
  {$-\dfrac{9}{2}$}{$\dfrac{9 \pi}{2}$}{$9$}

\cauTN[id=C4-B13-D02-NB-030,muc=NB,loigiai={Theo công thức: $S=\displaystyle\int_{a}^{b}\left|u(x)-v(x)\right|\,\mathrm{d}x$.}]{Cho các hàm số $y=u(x)$, $y=v(x)$ liên tục trên đoạn $[a;b]$. Diện tích $S$ của hình phẳng giới hạn bởi đồ thị hai hàm số đó và hai đường thẳng $x=a$, $x=b$ được tính theo công thức}
  {$S=\displaystyle\int_{a}^{b}\left|u(x)-v(x)\right|\,\mathrm{d}x$}
  {$S=\pi\displaystyle\int_{a}^{b}\left|u(x)-v(x)\right|\,\mathrm{d}x$}{$S=\displaystyle\int_{a}^{b}\left|u(x)+v(x)\right|\,\mathrm{d}x$}{$S=\pi\displaystyle\int_{a}^{b}\left[u^2(x)-v^2(x)\right]\,\mathrm{d}x$}

\cauTN[id=C4-B13-D02-NB-031,muc=NB,loigiai={Phương trình hoành độ giao điểm: $- x^{2} - 4 x + 3=3 - x\Leftrightarrow - x^{2} - 3 x=0\Leftrightarrow x\in\{-3;0\}$. $S=\displaystyle\int_{-3}^{0}\left|- x^{2} - 3 x\right|\,\mathrm{d}x=\dfrac{9}{2}$.}]{Diện tích hình phẳng giới hạn bởi parabol $y=- x^{2} - 4 x + 3$ và đường thẳng $y=3 - x$ bằng}
  {$\dfrac{9}{2}$}
  {$-\dfrac{9}{2}$}{$\dfrac{9 \pi}{2}$}{$9$}

\cauTN[id=C4-B13-D02-NB-032,muc=NB,loigiai={Phương trình hoành độ giao điểm: $x^{2} + x - 1=2 x + 1\Leftrightarrow x^{2} - x - 2=0\Leftrightarrow x\in\{-1;2\}$. $S=\displaystyle\int_{-1}^{2}\left|x^{2} - x - 2\right|\,\mathrm{d}x=\dfrac{9}{2}$.}]{Diện tích hình phẳng giới hạn bởi parabol $y=x^{2} + x - 1$ và đường thẳng $y=2 x + 1$ bằng}
  {$\dfrac{9}{2}$}
  {$-\dfrac{9}{2}$}{$\dfrac{9 \pi}{2}$}{$9$}

\cauTN[id=C4-B13-D02-NB-033,muc=NB,loigiai={Phương trình hoành độ giao điểm: $x^{2} - x - 2=- 2 x\Leftrightarrow x^{2} + x - 2=0\Leftrightarrow x\in\{-2;1\}$. $S=\displaystyle\int_{-2}^{1}\left|x^{2} + x - 2\right|\,\mathrm{d}x=\dfrac{9}{2}$.}]{Diện tích hình phẳng giới hạn bởi parabol $y=x^{2} - x - 2$ và đường thẳng $y=- 2 x$ bằng}
  {$\dfrac{9}{2}$}
  {$-\dfrac{9}{2}$}{$\dfrac{9 \pi}{2}$}{$9$}

\cauTN[id=C4-B13-D02-NB-034,muc=NB,loigiai={Phương trình hoành độ giao điểm: $x^{2} - x - 6=- x - 2\Leftrightarrow x^{2} - 4=0\Leftrightarrow x\in\{-2;2\}$. $S=\displaystyle\int_{-2}^{2}\left|x^{2} - 4\right|\,\mathrm{d}x=\dfrac{32}{3}$.}]{Diện tích hình phẳng giới hạn bởi parabol $y=x^{2} - x - 6$ và đường thẳng $y=- x - 2$ bằng}
  {$\dfrac{32}{3}$}
  {$-\dfrac{32}{3}$}{$\dfrac{32 \pi}{3}$}{$\dfrac{64}{3}$}

\cauTN[id=C4-B13-D02-NB-035,muc=NB,loigiai={Theo công thức: $S=\displaystyle\int_{0}^{2}\left|f_1(x)-f_2(x)\right|\,\mathrm{d}x$.}]{Cho các hàm số $y=f_1(x)$, $y=f_2(x)$ liên tục trên đoạn $[0;2]$. Diện tích $S$ của hình phẳng giới hạn bởi đồ thị hai hàm số đó và hai đường thẳng $x=0$, $x=2$ được tính theo công thức}
  {$S=\displaystyle\int_{0}^{2}\left|f_1(x)-f_2(x)\right|\,\mathrm{d}x$}
  {$S=\displaystyle\int_{0}^{2}\left|f_1(x)+f_2(x)\right|\,\mathrm{d}x$}{$S=\pi\displaystyle\int_{0}^{2}\left[f_1^2(x)-f_2^2(x)\right]\,\mathrm{d}x$}{$S=\displaystyle\int_{0}^{2}\left[f_1(x)-f_2(x)\right]\,\mathrm{d}x$}

\cauTN[id=C4-B13-D02-NB-036,muc=NB,loigiai={Phương trình hoành độ giao điểm: $- x^{2} - x - 3=2 x - 3\Leftrightarrow - x^{2} - 3 x=0\Leftrightarrow x\in\{-3;0\}$. $S=\displaystyle\int_{-3}^{0}\left|- x^{2} - 3 x\right|\,\mathrm{d}x=\dfrac{9}{2}$.}]{Diện tích hình phẳng giới hạn bởi parabol $y=- x^{2} - x - 3$ và đường thẳng $y=2 x - 3$ bằng}
  {$\dfrac{9}{2}$}
  {$-\dfrac{9}{2}$}{$\dfrac{9 \pi}{2}$}{$9$}

\cauDS[id=C4-B13-D02-NB-037,muc=NB,loigiai={Phương trình hoành độ giao điểm: $x^{2} - 2 x=1 - 2 x\Leftrightarrow x^{2} - 1=0\Leftrightarrow x\in\{-1;1\}$. $S=\displaystyle\int_{-1}^{1}\left|x^{2} - 1\right|\,\mathrm{d}x=\dfrac{4}{3}$. a) Đúng. b) Sai: không có thừa số $\pi$. c) Sai: phải lấy hiệu hai hàm số. d) Đúng.}]{Gọi $S$ là diện tích hình phẳng giới hạn bởi đồ thị các hàm số $y=x^{2} - 2 x$ và $y=1 - 2 x$. Xét tính đúng sai của các mệnh đề sau:}{%
  \y{$S=\displaystyle\int_{-1}^{1}\left|x^{2} - 1\right|\,\mathrm{d}x$.}{d}
  \y{$S=\pi\displaystyle\int_{-1}^{1}\left|x^{2} - 1\right|\,\mathrm{d}x$.}{s}
  \y{$S=\displaystyle\int_{-1}^{1}\left|x^{2} - 4 x + 1\right|\,\mathrm{d}x$.}{s}
  \y{$S=\dfrac{4}{3}$.}{d}}

\cauDS[id=C4-B13-D02-NB-038,muc=NB,loigiai={Phương trình hoành độ giao điểm: $3 x^{2} - 3 x - 9=x^{2} - x + 3\Leftrightarrow 2 x^{2} - 2 x - 12=0\Leftrightarrow x\in\{-2;3\}$. $S=\displaystyle\int_{-2}^{3}\left|2 x^{2} - 2 x - 12\right|\,\mathrm{d}x=\dfrac{125}{3}$. a) Sai: đúng là $\dfrac{125}{3}$. b) Đúng. c) Sai: không có thừa số $\pi$. d) Sai: phải lấy hiệu hai hàm số.}]{Gọi $S$ là diện tích hình phẳng giới hạn bởi đồ thị các hàm số $y=3 x^{2} - 3 x - 9$ và $y=x^{2} - x + 3$. Xét tính đúng sai của các mệnh đề sau:}{%
  \y{$S=\dfrac{250}{3}$.}{s}
  \y{Hai đồ thị cắt nhau tại hai điểm có hoành độ $x=-2$ và $x=3$.}{d}
  \y{$S=\pi\displaystyle\int_{-2}^{3}\left|2 x^{2} - 2 x - 12\right|\,\mathrm{d}x$.}{s}
  \y{$S=\displaystyle\int_{-2}^{3}\left|4 x^{2} - 4 x - 6\right|\,\mathrm{d}x$.}{s}}

\cauDS[id=C4-B13-D02-NB-039,muc=NB,loigiai={Phương trình hoành độ giao điểm: $x^{2} - 3 x + 3=x\Leftrightarrow x^{2} - 4 x + 3=0\Leftrightarrow x\in\{1;3\}$. $S=\displaystyle\int_{1}^{3}\left|x^{2} - 4 x + 3\right|\,\mathrm{d}x=\dfrac{4}{3}$. a) Đúng. b) Đúng. c) Sai: đúng là $\dfrac{4}{3}$. d) Đúng.}]{Gọi $S$ là diện tích hình phẳng giới hạn bởi đồ thị các hàm số $y=x^{2} - 3 x + 3$ và $y=x$. Xét tính đúng sai của các mệnh đề sau:}{%
  \y{$S=\dfrac{4}{3}$.}{d}
  \y{$S=\displaystyle\int_{1}^{3}\left|x^{2} - 4 x + 3\right|\,\mathrm{d}x$.}{d}
  \y{$S=\dfrac{8}{3}$.}{s}
  \y{Hai đồ thị cắt nhau tại hai điểm có hoành độ $x=1$ và $x=3$.}{d}}

\cauDS[id=C4-B13-D02-NB-040,muc=NB,loigiai={Phương trình hoành độ giao điểm: $- x^{2} + x + 5=-1\Leftrightarrow - x^{2} + x + 6=0\Leftrightarrow x\in\{-2;3\}$. $S=\displaystyle\int_{-2}^{3}\left|- x^{2} + x + 6\right|\,\mathrm{d}x=\dfrac{125}{6}$. a) Đúng. b) Sai: phải lấy hiệu hai hàm số. c) Sai: đúng là $\dfrac{125}{6}$. d) Sai: không có thừa số $\pi$.}]{Gọi $S$ là diện tích hình phẳng giới hạn bởi đồ thị các hàm số $y=- x^{2} + x + 5$ và $y=-1$. Xét tính đúng sai của các mệnh đề sau:}{%
  \y{$S=\displaystyle\int_{-2}^{3}\left|- x^{2} + x + 6\right|\,\mathrm{d}x$.}{d}
  \y{$S=\displaystyle\int_{-2}^{3}\left|- x^{2} + x + 4\right|\,\mathrm{d}x$.}{s}
  \y{$S=\dfrac{125}{3}$.}{s}
  \y{$S=\pi\displaystyle\int_{-2}^{3}\left|- x^{2} + x + 6\right|\,\mathrm{d}x$.}{s}}

\cauDS[id=C4-B13-D02-NB-041,muc=NB,loigiai={Phương trình hoành độ giao điểm: $x^{2} - 5 x + 1=- 2 x - 1\Leftrightarrow x^{2} - 3 x + 2=0\Leftrightarrow x\in\{1;2\}$. $S=\displaystyle\int_{1}^{2}\left|x^{2} - 3 x + 2\right|\,\mathrm{d}x=\dfrac{1}{6}$. a) Đúng. b) Sai: phải lấy hiệu hai hàm số. c) Sai: không có thừa số $\pi$. d) Sai: đúng là $\dfrac{1}{6}$.}]{Gọi $S$ là diện tích hình phẳng giới hạn bởi đồ thị các hàm số $y=x^{2} - 5 x + 1$ và $y=- 2 x - 1$. Xét tính đúng sai của các mệnh đề sau:}{%
  \y{$S=\displaystyle\int_{1}^{2}\left|x^{2} - 3 x + 2\right|\,\mathrm{d}x$.}{d}
  \y{$S=\displaystyle\int_{1}^{2}\left|x^{2} - 7 x\right|\,\mathrm{d}x$.}{s}
  \y{$S=\pi\displaystyle\int_{1}^{2}\left|x^{2} - 3 x + 2\right|\,\mathrm{d}x$.}{s}
  \y{$S=\dfrac{1}{3}$.}{s}}

\cauDS[id=C4-B13-D02-NB-042,muc=NB,loigiai={Phương trình hoành độ giao điểm: $- x^{2} - 3 x + 4=- 2 x - 2\Leftrightarrow - x^{2} - x + 6=0\Leftrightarrow x\in\{-3;2\}$. $S=\displaystyle\int_{-3}^{2}\left|- x^{2} - x + 6\right|\,\mathrm{d}x=\dfrac{125}{6}$. a) Đúng. b) Sai: phải lấy hiệu hai hàm số. c) Sai: không có thừa số $\pi$. d) Đúng.}]{Gọi $S$ là diện tích hình phẳng giới hạn bởi đồ thị các hàm số $y=- x^{2} - 3 x + 4$ và $y=- 2 x - 2$. Xét tính đúng sai của các mệnh đề sau:}{%
  \y{$S=\displaystyle\int_{-3}^{2}\left|- x^{2} - x + 6\right|\,\mathrm{d}x$.}{d}
  \y{$S=\displaystyle\int_{-3}^{2}\left|- x^{2} - 5 x + 2\right|\,\mathrm{d}x$.}{s}
  \y{$S=\pi\displaystyle\int_{-3}^{2}\left|- x^{2} - x + 6\right|\,\mathrm{d}x$.}{s}
  \y{$S=\dfrac{125}{6}$.}{d}}

\cauDS[id=C4-B13-D02-NB-043,muc=NB,loigiai={Phương trình hoành độ giao điểm: $- x^{2} - x + 4=x + 1\Leftrightarrow - x^{2} - 2 x + 3=0\Leftrightarrow x\in\{-3;1\}$. $S=\displaystyle\int_{-3}^{1}\left|- x^{2} - 2 x + 3\right|\,\mathrm{d}x=\dfrac{32}{3}$. a) Đúng. b) Sai: đúng là $\dfrac{32}{3}$. c) Đúng. d) Sai: phải lấy hiệu hai hàm số.}]{Gọi $S$ là diện tích hình phẳng giới hạn bởi đồ thị các hàm số $y=- x^{2} - x + 4$ và $y=x + 1$. Xét tính đúng sai của các mệnh đề sau:}{%
  \y{$S=\displaystyle\int_{-3}^{1}\left|- x^{2} - 2 x + 3\right|\,\mathrm{d}x$.}{d}
  \y{$S=\dfrac{64}{3}$.}{s}
  \y{Hai đồ thị cắt nhau tại hai điểm có hoành độ $x=-3$ và $x=1$.}{d}
  \y{$S=\displaystyle\int_{-3}^{1}\left|5 - x^{2}\right|\,\mathrm{d}x$.}{s}}

\cauDS[id=C4-B13-D02-NB-044,muc=NB,loigiai={Phương trình hoành độ giao điểm: $x^{2} - x - 8=-2\Leftrightarrow x^{2} - x - 6=0\Leftrightarrow x\in\{-2;3\}$. $S=\displaystyle\int_{-2}^{3}\left|x^{2} - x - 6\right|\,\mathrm{d}x=\dfrac{125}{6}$. a) Sai: đúng là $\dfrac{125}{6}$. b) Sai: phải lấy hiệu hai hàm số. c) Sai: không có thừa số $\pi$. d) Đúng.}]{Gọi $S$ là diện tích hình phẳng giới hạn bởi đồ thị các hàm số $y=x^{2} - x - 8$ và $y=-2$. Xét tính đúng sai của các mệnh đề sau:}{%
  \y{$S=\dfrac{125}{3}$.}{s}
  \y{$S=\displaystyle\int_{-2}^{3}\left|x^{2} - x - 10\right|\,\mathrm{d}x$.}{s}
  \y{$S=\pi\displaystyle\int_{-2}^{3}\left|x^{2} - x - 6\right|\,\mathrm{d}x$.}{s}
  \y{$S=\dfrac{125}{6}$.}{d}}

\cauDS[id=C4-B13-D02-NB-045,muc=NB,loigiai={Phương trình hoành độ giao điểm: $3 - x^{2}=-1\Leftrightarrow 4 - x^{2}=0\Leftrightarrow x\in\{-2;2\}$. $S=\displaystyle\int_{-2}^{2}\left|4 - x^{2}\right|\,\mathrm{d}x=\dfrac{32}{3}$. a) Đúng. b) Sai: đúng là $\dfrac{32}{3}$. c) Đúng. d) Đúng.}]{Gọi $S$ là diện tích hình phẳng giới hạn bởi đồ thị các hàm số $y=3 - x^{2}$ và $y=-1$. Xét tính đúng sai của các mệnh đề sau:}{%
  \y{$S=\dfrac{32}{3}$.}{d}
  \y{$S=\dfrac{64}{3}$.}{s}
  \y{$S=\displaystyle\int_{-2}^{2}\left|4 - x^{2}\right|\,\mathrm{d}x$.}{d}
  \y{Hai đồ thị cắt nhau tại hai điểm có hoành độ $x=-2$ và $x=2$.}{d}}

\cauDS[id=C4-B13-D02-NB-046,muc=NB,loigiai={Phương trình hoành độ giao điểm: $x^{2} - 2 x - 1=- x^{2} - 1\Leftrightarrow 2 x^{2} - 2 x=0\Leftrightarrow x\in\{0;1\}$. $S=\displaystyle\int_{0}^{1}\left|2 x^{2} - 2 x\right|\,\mathrm{d}x=\dfrac{1}{3}$. a) Đúng. b) Sai: không có thừa số $\pi$. c) Đúng. d) Sai: đúng là $\dfrac{1}{3}$.}]{Gọi $S$ là diện tích hình phẳng giới hạn bởi đồ thị các hàm số $y=x^{2} - 2 x - 1$ và $y=- x^{2} - 1$. Xét tính đúng sai của các mệnh đề sau:}{%
  \y{$S=\dfrac{1}{3}$.}{d}
  \y{$S=\pi\displaystyle\int_{0}^{1}\left|2 x^{2} - 2 x\right|\,\mathrm{d}x$.}{s}
  \y{$S=\displaystyle\int_{0}^{1}\left|2 x^{2} - 2 x\right|\,\mathrm{d}x$.}{d}
  \y{$S=\dfrac{2}{3}$.}{s}}

\cauDS[id=C4-B13-D02-NB-047,muc=NB,loigiai={Phương trình hoành độ giao điểm: $2 x^{2} + x - 1=x^{2} + x + 3\Leftrightarrow x^{2} - 4=0\Leftrightarrow x\in\{-2;2\}$. $S=\displaystyle\int_{-2}^{2}\left|x^{2} - 4\right|\,\mathrm{d}x=\dfrac{32}{3}$. a) Sai: không có thừa số $\pi$. b) Đúng. c) Đúng. d) Đúng.}]{Gọi $S$ là diện tích hình phẳng giới hạn bởi đồ thị các hàm số $y=2 x^{2} + x - 1$ và $y=x^{2} + x + 3$. Xét tính đúng sai của các mệnh đề sau:}{%
  \y{$S=\pi\displaystyle\int_{-2}^{2}\left|x^{2} - 4\right|\,\mathrm{d}x$.}{s}
  \y{$S=\displaystyle\int_{-2}^{2}\left|x^{2} - 4\right|\,\mathrm{d}x$.}{d}
  \y{$S=\dfrac{32}{3}$.}{d}
  \y{Hai đồ thị cắt nhau tại hai điểm có hoành độ $x=-2$ và $x=2$.}{d}}

\cauDS[id=C4-B13-D02-NB-048,muc=NB,loigiai={Phương trình hoành độ giao điểm: $- 2 x^{2} + 2 x - 2=- x^{2} + 2 x - 3\Leftrightarrow 1 - x^{2}=0\Leftrightarrow x\in\{-1;1\}$. $S=\displaystyle\int_{-1}^{1}\left|1 - x^{2}\right|\,\mathrm{d}x=\dfrac{4}{3}$. a) Sai: không có thừa số $\pi$. b) Đúng. c) Đúng. d) Đúng.}]{Gọi $S$ là diện tích hình phẳng giới hạn bởi đồ thị các hàm số $y=- 2 x^{2} + 2 x - 2$ và $y=- x^{2} + 2 x - 3$. Xét tính đúng sai của các mệnh đề sau:}{%
  \y{$S=\pi\displaystyle\int_{-1}^{1}\left|1 - x^{2}\right|\,\mathrm{d}x$.}{s}
  \y{$S=\dfrac{4}{3}$.}{d}
  \y{Hai đồ thị cắt nhau tại hai điểm có hoành độ $x=-1$ và $x=1$.}{d}
  \y{$S=\displaystyle\int_{-1}^{1}\left|1 - x^{2}\right|\,\mathrm{d}x$.}{d}}

\cauTLN[id=C4-B13-D02-NB-049,muc=NB,loigiai={Phương trình hoành độ giao điểm: $- x^{2} + x + 9=2 x + 3\Leftrightarrow - x^{2} - x + 6=0\Leftrightarrow x\in\{-3;2\}$. $S=\displaystyle\int_{-3}^{2}\left|- x^{2} - x + 6\right|\,\mathrm{d}x=\dfrac{125}{6}$.}]{Tính diện tích hình phẳng giới hạn bởi đồ thị các hàm số $y=- x^{2} + x + 9$ và $y=2 x + 3$ (làm tròn kết quả đến hàng phần mười).}{20,8}

\cauTLN[id=C4-B13-D02-NB-050,muc=NB,loigiai={Phương trình hoành độ giao điểm: $2 x^{2} + x - 1=x^{2} + x\Leftrightarrow x^{2} - 1=0\Leftrightarrow x\in\{-1;1\}$. $S=\displaystyle\int_{-1}^{1}\left|x^{2} - 1\right|\,\mathrm{d}x=\dfrac{4}{3}$.}]{Tính diện tích hình phẳng giới hạn bởi đồ thị các hàm số $y=2 x^{2} + x - 1$ và $y=x^{2} + x$ (làm tròn kết quả đến hàng phần trăm).}{1,33}

\cauTLN[id=C4-B13-D02-NB-051,muc=NB,loigiai={Phương trình hoành độ giao điểm: $x^{2} + x + 1=2 x + 3\Leftrightarrow x^{2} - x - 2=0\Leftrightarrow x\in\{-1;2\}$. $S=\displaystyle\int_{-1}^{2}\left|x^{2} - x - 2\right|\,\mathrm{d}x=\dfrac{9}{2}$.}]{Tính diện tích hình phẳng giới hạn bởi đồ thị các hàm số $y=x^{2} + x + 1$ và $y=2 x + 3$.}{4,5}

\cauTLN[id=C4-B13-D02-NB-052,muc=NB,loigiai={Phương trình hoành độ giao điểm: $- 2 x^{2} - 3 x + 4=- x^{2} - 2 x - 2\Leftrightarrow - x^{2} - x + 6=0\Leftrightarrow x\in\{-3;2\}$. $S=\displaystyle\int_{-3}^{2}\left|- x^{2} - x + 6\right|\,\mathrm{d}x=\dfrac{125}{6}$.}]{Tính diện tích hình phẳng giới hạn bởi đồ thị các hàm số $y=- 2 x^{2} - 3 x + 4$ và $y=- x^{2} - 2 x - 2$ (làm tròn kết quả đến hàng phần mười).}{20,8}

\cauTLN[id=C4-B13-D02-NB-053,muc=NB,loigiai={Phương trình hoành độ giao điểm: $- x^{2} + x + 5=2 x - 1\Leftrightarrow - x^{2} - x + 6=0\Leftrightarrow x\in\{-3;2\}$. $S=\displaystyle\int_{-3}^{2}\left|- x^{2} - x + 6\right|\,\mathrm{d}x=\dfrac{125}{6}$.}]{Tính diện tích hình phẳng giới hạn bởi đồ thị các hàm số $y=- x^{2} + x + 5$ và $y=2 x - 1$ (làm tròn kết quả đến hàng phần mười).}{20,8}

\cauTLN[id=C4-B13-D02-NB-054,muc=NB,loigiai={Phương trình hoành độ giao điểm: $- x^{2} + 3 x + 2=2 x + 2\Leftrightarrow - x^{2} + x=0\Leftrightarrow x\in\{0;1\}$. $S=\displaystyle\int_{0}^{1}\left|- x^{2} + x\right|\,\mathrm{d}x=\dfrac{1}{6}$.}]{Tính diện tích hình phẳng giới hạn bởi đồ thị các hàm số $y=- x^{2} + 3 x + 2$ và $y=2 x + 2$ (làm tròn kết quả đến hàng phần trăm).}{0,17}

\cauTLN[id=C4-B13-D02-NB-055,muc=NB,loigiai={Phương trình hoành độ giao điểm: $2 - x^{2}=2 - 2 x\Leftrightarrow - x^{2} + 2 x=0\Leftrightarrow x\in\{0;2\}$. $S=\displaystyle\int_{0}^{2}\left|- x^{2} + 2 x\right|\,\mathrm{d}x=\dfrac{4}{3}$.}]{Tính diện tích hình phẳng giới hạn bởi đồ thị các hàm số $y=2 - x^{2}$ và $y=2 - 2 x$ (làm tròn kết quả đến hàng phần trăm).}{1,33}

\cauTLN[id=C4-B13-D02-NB-056,muc=NB,loigiai={Phương trình hoành độ giao điểm: $x^{2} - 3 x + 8=2 x + 2\Leftrightarrow x^{2} - 5 x + 6=0\Leftrightarrow x\in\{2;3\}$. $S=\displaystyle\int_{2}^{3}\left|x^{2} - 5 x + 6\right|\,\mathrm{d}x=\dfrac{1}{6}$.}]{Tính diện tích hình phẳng giới hạn bởi đồ thị các hàm số $y=x^{2} - 3 x + 8$ và $y=2 x + 2$ (làm tròn kết quả đến hàng phần trăm).}{0,17}

\cauTLN[id=C4-B13-D02-NB-057,muc=NB,loigiai={Phương trình hoành độ giao điểm: $x^{2} - 2 x - 2=-2\Leftrightarrow x^{2} - 2 x=0\Leftrightarrow x\in\{0;2\}$. $S=\displaystyle\int_{0}^{2}\left|x^{2} - 2 x\right|\,\mathrm{d}x=\dfrac{4}{3}$.}]{Tính diện tích hình phẳng giới hạn bởi đồ thị các hàm số $y=x^{2} - 2 x - 2$ và $y=-2$ (làm tròn kết quả đến hàng phần trăm).}{1,33}

\cauTLN[id=C4-B13-D02-NB-058,muc=NB,loigiai={Phương trình hoành độ giao điểm: $x^{2} - 2 x + 1=1 - x\Leftrightarrow x^{2} - x=0\Leftrightarrow x\in\{0;1\}$. $S=\displaystyle\int_{0}^{1}\left|x^{2} - x\right|\,\mathrm{d}x=\dfrac{1}{6}$.}]{Tính diện tích hình phẳng giới hạn bởi đồ thị các hàm số $y=x^{2} - 2 x + 1$ và $y=1 - x$ (làm tròn kết quả đến hàng phần trăm).}{0,17}

\cauTLN[id=C4-B13-D02-NB-059,muc=NB,loigiai={Phương trình hoành độ giao điểm: $- x^{2} + 2 x - 2=1 - 2 x\Leftrightarrow - x^{2} + 4 x - 3=0\Leftrightarrow x\in\{1;3\}$. $S=\displaystyle\int_{1}^{3}\left|- x^{2} + 4 x - 3\right|\,\mathrm{d}x=\dfrac{4}{3}$.}]{Tính diện tích hình phẳng giới hạn bởi đồ thị các hàm số $y=- x^{2} + 2 x - 2$ và $y=1 - 2 x$ (làm tròn kết quả đến hàng phần trăm).}{1,33}

\cauTLN[id=C4-B13-D02-NB-060,muc=NB,loigiai={Phương trình hoành độ giao điểm: $- x^{2} + x + 3=2 x + 3\Leftrightarrow - x^{2} - x=0\Leftrightarrow x\in\{-1;0\}$. $S=\displaystyle\int_{-1}^{0}\left|- x^{2} - x\right|\,\mathrm{d}x=\dfrac{1}{6}$.}]{Tính diện tích hình phẳng giới hạn bởi đồ thị các hàm số $y=- x^{2} + x + 3$ và $y=2 x + 3$ (làm tròn kết quả đến hàng phần trăm).}{0,17}

\cauTN[id=C4-B13-D02-TH-001,muc=TH,loigiai={$S=\displaystyle\int_{0}^{1}\left|x^{2} + 3-0\right|\,\mathrm{d}x=\dfrac{10}{3}$.}]{Diện tích hình phẳng giới hạn bởi đồ thị các hàm số $y=x^{2} + 3$, $y=0$ và hai đường thẳng $x=0$, $x=1$ bằng}
  {$\dfrac{10}{3}$}
  {$\dfrac{10 \pi}{3}$}{$\dfrac{13}{3}$}{$\dfrac{20}{3}$}

\cauTN[id=C4-B13-D02-TH-002,muc=TH,loigiai={Phương trình hoành độ giao điểm: $- x^{2} + 3 x + 4=2 x + 2\Leftrightarrow - x^{2} + x + 2=0\Leftrightarrow x\in\{-1;2\}$. $S=\displaystyle\int_{-1}^{2}\left|- x^{2} + x + 2\right|\,\mathrm{d}x=\dfrac{9}{2}$.}]{Diện tích hình phẳng giới hạn bởi đồ thị các hàm số $y=- x^{2} + 3 x + 4$ và $y=2 x + 2$ bằng}
  {$\dfrac{9}{2}$}
  {$-\dfrac{9}{2}$}{$\dfrac{9 \pi}{2}$}{$9$}

\cauTN[id=C4-B13-D02-TH-003,muc=TH,loigiai={$S=\displaystyle\int_{\dfrac{\pi}{4}}^{\pi}\left|\sin x-\cos x\right|\,\mathrm{d}x=1 + \sqrt{2}$.}]{Diện tích hình phẳng giới hạn bởi đồ thị các hàm số $y=\sin x$, $y=\cos x$ và hai đường thẳng $x=\dfrac{\pi}{4}$, $x=\pi$ bằng}
  {$1 + \sqrt{2}$}
  {$1$}{$\dfrac{\sqrt{2}}{2} + 1$}{$\pi (1 + \sqrt{2})$}

\cauTN[id=C4-B13-D02-TH-004,muc=TH,loigiai={$S=\displaystyle\int_{0}^{\dfrac{\pi}{2}}\left|\sin x-\cos x\right|\,\mathrm{d}x=-2 + 2 \sqrt{2}$.}]{Diện tích hình phẳng giới hạn bởi đồ thị các hàm số $y=\sin x$, $y=\cos x$ và hai đường thẳng $x=0$, $x=\dfrac{\pi}{2}$ bằng}
  {$-2 + 2 \sqrt{2}$}
  {$2$}{$1$}{$2 \pi (-1 + \sqrt{2})$}

\cauTN[id=C4-B13-D02-TH-005,muc=TH,loigiai={$S=\displaystyle\int_{0}^{2}\left|x^{2} + 3-x\right|\,\mathrm{d}x=\dfrac{20}{3}$.}]{Diện tích hình phẳng giới hạn bởi đồ thị các hàm số $y=x^{2} + 3$, $y=x$ và hai đường thẳng $x=0$, $x=2$ bằng}
  {$\dfrac{20}{3}$}
  {$\dfrac{32}{3}$}{$\dfrac{26}{3}$}{$\dfrac{20 \pi}{3}$}

\cauTN[id=C4-B13-D02-TH-006,muc=TH,loigiai={$S=\displaystyle\int_{0}^{2}\left|x^{2} + 3-0\right|\,\mathrm{d}x=\dfrac{26}{3}$.}]{Diện tích hình phẳng giới hạn bởi đồ thị các hàm số $y=x^{2} + 3$, $y=0$ và hai đường thẳng $x=0$, $x=2$ bằng}
  {$\dfrac{26}{3}$}
  {$\dfrac{26 \pi}{3}$}{$\dfrac{29}{3}$}{$\dfrac{52}{3}$}

\cauTN[id=C4-B13-D02-TH-007,muc=TH,loigiai={Phương trình hoành độ giao điểm: $- x^{2} + x - 2=2 x - 2\Leftrightarrow - x^{2} - x=0\Leftrightarrow x\in\{-1;0\}$. $S=\displaystyle\int_{-1}^{0}\left|- x^{2} - x\right|\,\mathrm{d}x=\dfrac{1}{6}$.}]{Diện tích hình phẳng giới hạn bởi đồ thị các hàm số $y=- x^{2} + x - 2$ và $y=2 x - 2$ bằng}
  {$\dfrac{1}{6}$}
  {$-\dfrac{1}{6}$}{$\dfrac{\pi}{6}$}{$\dfrac{1}{3}$}

\cauTN[id=C4-B13-D02-TH-008,muc=TH,loigiai={Phương trình hoành độ giao điểm: $2 x^{2} - 3 x - 2=x^{2} - 2 x\Leftrightarrow x^{2} - x - 2=0\Leftrightarrow x\in\{-1;2\}$. $S=\displaystyle\int_{-1}^{2}\left|x^{2} - x - 2\right|\,\mathrm{d}x=\dfrac{9}{2}$.}]{Diện tích hình phẳng giới hạn bởi đồ thị các hàm số $y=2 x^{2} - 3 x - 2$ và $y=x^{2} - 2 x$ bằng}
  {$\dfrac{9}{2}$}
  {$\dfrac{9 \pi}{2}$}{$9$}{$\dfrac{11}{2}$}

\cauTN[id=C4-B13-D02-TH-009,muc=TH,loigiai={Phương trình hoành độ giao điểm: $x^{2} + 4 x + 6=3\Leftrightarrow x^{2} + 4 x + 3=0\Leftrightarrow x\in\{-3;-1\}$. $S=\displaystyle\int_{-3}^{-1}\left|x^{2} + 4 x + 3\right|\,\mathrm{d}x=\dfrac{4}{3}$.}]{Diện tích hình phẳng giới hạn bởi đồ thị các hàm số $y=x^{2} + 4 x + 6$ và $y=3$ bằng}
  {$\dfrac{4}{3}$}
  {$\dfrac{4 \pi}{3}$}{$\dfrac{8}{3}$}{$\dfrac{7}{3}$}

\cauTN[id=C4-B13-D02-TH-010,muc=TH,loigiai={Phương trình hoành độ giao điểm: $3 x^{2} + x - 5=x^{2} - x - 1\Leftrightarrow 2 x^{2} + 2 x - 4=0\Leftrightarrow x\in\{-2;1\}$. $S=\displaystyle\int_{-2}^{1}\left|2 x^{2} + 2 x - 4\right|\,\mathrm{d}x=9$.}]{Diện tích hình phẳng giới hạn bởi đồ thị các hàm số $y=3 x^{2} + x - 5$ và $y=x^{2} - x - 1$ bằng}
  {$9$}
  {$9 \pi$}{$18$}{$10$}

\cauTN[id=C4-B13-D02-TH-011,muc=TH,loigiai={Phương trình hoành độ giao điểm: $2 x^{2} - 2=x^{2} + 2 x + 1\Leftrightarrow x^{2} - 2 x - 3=0\Leftrightarrow x\in\{-1;3\}$. $S=\displaystyle\int_{-1}^{3}\left|x^{2} - 2 x - 3\right|\,\mathrm{d}x=\dfrac{32}{3}$.}]{Diện tích hình phẳng giới hạn bởi đồ thị các hàm số $y=2 x^{2} - 2$ và $y=x^{2} + 2 x + 1$ bằng}
  {$\dfrac{32}{3}$}
  {$\dfrac{32 \pi}{3}$}{$\dfrac{64}{3}$}{$\dfrac{35}{3}$}

\cauTN[id=C4-B13-D02-TH-012,muc=TH,loigiai={Phương trình hoành độ giao điểm: $- 2 x^{2} + 6 x - 2=- x^{2} + 2 x + 1\Leftrightarrow - x^{2} + 4 x - 3=0\Leftrightarrow x\in\{1;3\}$. $S=\displaystyle\int_{1}^{3}\left|- x^{2} + 4 x - 3\right|\,\mathrm{d}x=\dfrac{4}{3}$.}]{Diện tích hình phẳng giới hạn bởi đồ thị các hàm số $y=- 2 x^{2} + 6 x - 2$ và $y=- x^{2} + 2 x + 1$ bằng}
  {$\dfrac{4}{3}$}
  {$-\dfrac{4}{3}$}{$\dfrac{4 \pi}{3}$}{$\dfrac{8}{3}$}

\cauTN[id=C4-B13-D02-TH-013,muc=TH,loigiai={Phương trình hoành độ giao điểm: $x^{2} - 3 x - 5=- 2 x - 3\Leftrightarrow x^{2} - x - 2=0\Leftrightarrow x\in\{-1;2\}$. $S=\displaystyle\int_{-1}^{2}\left|x^{2} - x - 2\right|\,\mathrm{d}x=\dfrac{9}{2}$.}]{Diện tích hình phẳng giới hạn bởi đồ thị các hàm số $y=x^{2} - 3 x - 5$ và $y=- 2 x - 3$ bằng}
  {$\dfrac{9}{2}$}
  {$\dfrac{9 \pi}{2}$}{$9$}{$\dfrac{11}{2}$}

\cauTN[id=C4-B13-D02-TH-014,muc=TH,loigiai={Phương trình hoành độ giao điểm: $x^{2} + 4 x + 1=-2\Leftrightarrow x^{2} + 4 x + 3=0\Leftrightarrow x\in\{-3;-1\}$. $S=\displaystyle\int_{-3}^{-1}\left|x^{2} + 4 x + 3\right|\,\mathrm{d}x=\dfrac{4}{3}$.}]{Diện tích hình phẳng giới hạn bởi đồ thị các hàm số $y=x^{2} + 4 x + 1$ và $y=-2$ bằng}
  {$\dfrac{4}{3}$}
  {$\dfrac{4 \pi}{3}$}{$\dfrac{8}{3}$}{$\dfrac{7}{3}$}

\cauTN[id=C4-B13-D02-TH-015,muc=TH,loigiai={$S=\displaystyle\int_{0}^{\pi}\left|\sin x-\cos x\right|\,\mathrm{d}x=2 \sqrt{2}$.}]{Diện tích hình phẳng giới hạn bởi đồ thị các hàm số $y=\sin x$, $y=\cos x$ và hai đường thẳng $x=0$, $x=\pi$ bằng}
  {$2 \sqrt{2}$}
  {$2$}{$2 \sqrt{2} \pi$}{$4 \sqrt{2}$}

\cauTN[id=C4-B13-D02-TH-016,muc=TH,loigiai={$S=\displaystyle\int_{1}^{3}\left|x^{2} + 2-0\right|\,\mathrm{d}x=\dfrac{38}{3}$.}]{Diện tích hình phẳng giới hạn bởi đồ thị các hàm số $y=x^{2} + 2$, $y=0$ và hai đường thẳng $x=1$, $x=3$ bằng}
  {$\dfrac{38}{3}$}
  {$\dfrac{38 \pi}{3}$}{$\dfrac{41}{3}$}{$\dfrac{76}{3}$}

\cauTN[id=C4-B13-D02-TH-017,muc=TH,loigiai={Phương trình hoành độ giao điểm: $x^{2} + 7 x + 7=2 x + 1\Leftrightarrow x^{2} + 5 x + 6=0\Leftrightarrow x\in\{-3;-2\}$. $S=\displaystyle\int_{-3}^{-2}\left|x^{2} + 5 x + 6\right|\,\mathrm{d}x=\dfrac{1}{6}$.}]{Diện tích hình phẳng giới hạn bởi đồ thị các hàm số $y=x^{2} + 7 x + 7$ và $y=2 x + 1$ bằng}
  {$\dfrac{1}{6}$}
  {$\dfrac{\pi}{6}$}{$\dfrac{1}{3}$}{$\dfrac{7}{6}$}

\cauTN[id=C4-B13-D02-TH-018,muc=TH,loigiai={$S=\displaystyle\int_{-1}^{1}\left|\mathrm{e}^x-x\right|\,\mathrm{d}x=\mathrm{e} - \dfrac{1}{\mathrm{e}}$.}]{Diện tích hình phẳng giới hạn bởi đồ thị các hàm số $y=\mathrm{e}^x$, $y=x$ và hai đường thẳng $x=-1$, $x=1$ bằng}
  {$\mathrm{e} - \dfrac{1}{\mathrm{e}}$}
  {$\pi (\mathrm{e} - \dfrac{1}{\mathrm{e}})$}{$- \dfrac{1}{\mathrm{e}} + 1 + \mathrm{e}$}{$- \dfrac{2}{\mathrm{e}} + 2 \mathrm{e}$}

\cauTN[id=C4-B13-D02-TH-019,muc=TH,loigiai={$S=\displaystyle\int_{0}^{1}\left|\mathrm{e}^x-x\right|\,\mathrm{d}x=- \dfrac{3}{2} + \mathrm{e}$.}]{Diện tích hình phẳng giới hạn bởi đồ thị các hàm số $y=\mathrm{e}^x$, $y=x$ và hai đường thẳng $x=0$, $x=1$ bằng}
  {$- \dfrac{3}{2} + \mathrm{e}$}
  {$- \dfrac{1}{2} + \mathrm{e}$}{$\mathrm{e} - 1$}{$\pi (- \dfrac{3}{2} + \mathrm{e})$}

\cauTN[id=C4-B13-D02-TH-020,muc=TH,loigiai={Phương trình hoành độ giao điểm: $x^{2} - x - 11=- x - 2\Leftrightarrow x^{2} - 9=0\Leftrightarrow x\in\{-3;3\}$. $S=\displaystyle\int_{-3}^{3}\left|x^{2} - 9\right|\,\mathrm{d}x=36$.}]{Diện tích hình phẳng giới hạn bởi đồ thị các hàm số $y=x^{2} - x - 11$ và $y=- x - 2$ bằng}
  {$36$}
  {$36 \pi$}{$72$}{$37$}

\cauTN[id=C4-B13-D02-TH-021,muc=TH,loigiai={Phương trình hoành độ giao điểm: $x^{2} - x - 2=- x - 1\Leftrightarrow x^{2} - 1=0\Leftrightarrow x\in\{-1;1\}$. $S=\displaystyle\int_{-1}^{1}\left|x^{2} - 1\right|\,\mathrm{d}x=\dfrac{4}{3}$.}]{Diện tích hình phẳng giới hạn bởi đồ thị các hàm số $y=x^{2} - x - 2$ và $y=- x - 1$ bằng}
  {$\dfrac{4}{3}$}
  {$\dfrac{4 \pi}{3}$}{$\dfrac{8}{3}$}{$\dfrac{7}{3}$}

\cauTN[id=C4-B13-D02-TH-022,muc=TH,loigiai={Phương trình hoành độ giao điểm: $7 - x^{2}=-2\Leftrightarrow 9 - x^{2}=0\Leftrightarrow x\in\{-3;3\}$. $S=\displaystyle\int_{-3}^{3}\left|9 - x^{2}\right|\,\mathrm{d}x=36$.}]{Diện tích hình phẳng giới hạn bởi đồ thị các hàm số $y=7 - x^{2}$ và $y=-2$ bằng}
  {$36$}
  {$-36$}{$36 \pi$}{$72$}

\cauTN[id=C4-B13-D02-TH-023,muc=TH,loigiai={$S=\displaystyle\int_{-1}^{1}\left|x^{2} + 2-(- x)\right|\,\mathrm{d}x=\dfrac{14}{3}$.}]{Diện tích hình phẳng giới hạn bởi đồ thị các hàm số $y=x^{2} + 2$, $y=- x$ và hai đường thẳng $x=-1$, $x=1$ bằng}
  {$\dfrac{14}{3}$}
  {$\dfrac{14 \pi}{3}$}{$\dfrac{17}{3}$}{$\dfrac{28}{3}$}

\cauTN[id=C4-B13-D02-TH-024,muc=TH,loigiai={Phương trình hoành độ giao điểm: $- 2 x^{2} + x + 1=1 - x^{2}\Leftrightarrow - x^{2} + x=0\Leftrightarrow x\in\{0;1\}$. $S=\displaystyle\int_{0}^{1}\left|- x^{2} + x\right|\,\mathrm{d}x=\dfrac{1}{6}$.}]{Diện tích hình phẳng giới hạn bởi đồ thị các hàm số $y=- 2 x^{2} + x + 1$ và $y=1 - x^{2}$ bằng}
  {$\dfrac{1}{6}$}
  {$-\dfrac{1}{6}$}{$\dfrac{\pi}{6}$}{$\dfrac{1}{3}$}

\cauTN[id=C4-B13-D02-TH-025,muc=TH,loigiai={Phương trình hoành độ giao điểm: $6 - x^{2}=3 - 2 x\Leftrightarrow - x^{2} + 2 x + 3=0\Leftrightarrow x\in\{-1;3\}$. $S=\displaystyle\int_{-1}^{3}\left|- x^{2} + 2 x + 3\right|\,\mathrm{d}x=\dfrac{32}{3}$.}]{Diện tích hình phẳng giới hạn bởi đồ thị các hàm số $y=6 - x^{2}$ và $y=3 - 2 x$ bằng}
  {$\dfrac{32}{3}$}
  {$-\dfrac{32}{3}$}{$\dfrac{32 \pi}{3}$}{$\dfrac{64}{3}$}

\cauTN[id=C4-B13-D02-TH-026,muc=TH,loigiai={$S=\displaystyle\int_{0}^{1}\left|\mathrm{e}^x-(- x)\right|\,\mathrm{d}x=- \dfrac{1}{2} + \mathrm{e}$.}]{Diện tích hình phẳng giới hạn bởi đồ thị các hàm số $y=\mathrm{e}^x$, $y=- x$ và hai đường thẳng $x=0$, $x=1$ bằng}
  {$- \dfrac{1}{2} + \mathrm{e}$}
  {$- \dfrac{3}{2} + \mathrm{e}$}{$\mathrm{e} - 1$}{$\pi (- \dfrac{1}{2} + \mathrm{e})$}

\cauTN[id=C4-B13-D02-TH-027,muc=TH,loigiai={Phương trình hoành độ giao điểm: $x^{2} + 2 x - 4=2 x - 3\Leftrightarrow x^{2} - 1=0\Leftrightarrow x\in\{-1;1\}$. $S=\displaystyle\int_{-1}^{1}\left|x^{2} - 1\right|\,\mathrm{d}x=\dfrac{4}{3}$.}]{Diện tích hình phẳng giới hạn bởi đồ thị các hàm số $y=x^{2} + 2 x - 4$ và $y=2 x - 3$ bằng}
  {$\dfrac{4}{3}$}
  {$\dfrac{4 \pi}{3}$}{$\dfrac{8}{3}$}{$\dfrac{7}{3}$}

\cauTN[id=C4-B13-D02-TH-028,muc=TH,loigiai={$S=\displaystyle\int_{1}^{2}\left|\mathrm{e}^x-x\right|\,\mathrm{d}x=- \mathrm{e} - \dfrac{3}{2} + \mathrm{e}^{2}$.}]{Diện tích hình phẳng giới hạn bởi đồ thị các hàm số $y=\mathrm{e}^x$, $y=x$ và hai đường thẳng $x=1$, $x=2$ bằng}
  {$- \mathrm{e} - \dfrac{3}{2} + \mathrm{e}^{2}$}
  {$- \mathrm{e} + \dfrac{3}{2} + \mathrm{e}^{2}$}{$- \mathrm{e} + \mathrm{e}^{2}$}{$\pi (- \mathrm{e} - \dfrac{3}{2} + \mathrm{e}^{2})$}

\cauTN[id=C4-B13-D02-TH-029,muc=TH,loigiai={Phương trình hoành độ giao điểm: $3 x^{2} - 4 x - 1=x^{2} - 2 x - 1\Leftrightarrow 2 x^{2} - 2 x=0\Leftrightarrow x\in\{0;1\}$. $S=\displaystyle\int_{0}^{1}\left|2 x^{2} - 2 x\right|\,\mathrm{d}x=\dfrac{1}{3}$.}]{Diện tích hình phẳng giới hạn bởi đồ thị các hàm số $y=3 x^{2} - 4 x - 1$ và $y=x^{2} - 2 x - 1$ bằng}
  {$\dfrac{1}{3}$}
  {$\dfrac{\pi}{3}$}{$\dfrac{2}{3}$}{$\dfrac{4}{3}$}

\cauTN[id=C4-B13-D02-TH-030,muc=TH,loigiai={Phương trình hoành độ giao điểm: $x^{2} - x - 3=3\Leftrightarrow x^{2} - x - 6=0\Leftrightarrow x\in\{-2;3\}$. $S=\displaystyle\int_{-2}^{3}\left|x^{2} - x - 6\right|\,\mathrm{d}x=\dfrac{125}{6}$.}]{Diện tích hình phẳng giới hạn bởi đồ thị các hàm số $y=x^{2} - x - 3$ và $y=3$ bằng}
  {$\dfrac{125}{6}$}
  {$\dfrac{125 \pi}{6}$}{$\dfrac{125}{3}$}{$\dfrac{131}{6}$}

\cauTN[id=C4-B13-D02-TH-031,muc=TH,loigiai={Phương trình hoành độ giao điểm: $- 2 x^{2} - 3 x - 2=- x^{2} - 2\Leftrightarrow - x^{2} - 3 x=0\Leftrightarrow x\in\{-3;0\}$. $S=\displaystyle\int_{-3}^{0}\left|- x^{2} - 3 x\right|\,\mathrm{d}x=\dfrac{9}{2}$.}]{Diện tích hình phẳng giới hạn bởi đồ thị các hàm số $y=- 2 x^{2} - 3 x - 2$ và $y=- x^{2} - 2$ bằng}
  {$\dfrac{9}{2}$}
  {$-\dfrac{9}{2}$}{$\dfrac{9 \pi}{2}$}{$9$}

\cauTN[id=C4-B13-D02-TH-032,muc=TH,loigiai={Phương trình hoành độ giao điểm: $2 x^{2} + 2 x + 1=x^{2} + 1\Leftrightarrow x^{2} + 2 x=0\Leftrightarrow x\in\{-2;0\}$. $S=\displaystyle\int_{-2}^{0}\left|x^{2} + 2 x\right|\,\mathrm{d}x=\dfrac{4}{3}$.}]{Diện tích hình phẳng giới hạn bởi đồ thị các hàm số $y=2 x^{2} + 2 x + 1$ và $y=x^{2} + 1$ bằng}
  {$\dfrac{4}{3}$}
  {$\dfrac{4 \pi}{3}$}{$\dfrac{8}{3}$}{$\dfrac{7}{3}$}

\cauTN[id=C4-B13-D02-TH-033,muc=TH,loigiai={Phương trình hoành độ giao điểm: $3 - x^{2}=2 x + 3\Leftrightarrow - x^{2} - 2 x=0\Leftrightarrow x\in\{-2;0\}$. $S=\displaystyle\int_{-2}^{0}\left|- x^{2} - 2 x\right|\,\mathrm{d}x=\dfrac{4}{3}$.}]{Diện tích hình phẳng giới hạn bởi đồ thị các hàm số $y=3 - x^{2}$ và $y=2 x + 3$ bằng}
  {$\dfrac{4}{3}$}
  {$-\dfrac{4}{3}$}{$\dfrac{4 \pi}{3}$}{$\dfrac{8}{3}$}

\cauTN[id=C4-B13-D02-TH-034,muc=TH,loigiai={Phương trình hoành độ giao điểm: $- x^{2}=- 2 x\Leftrightarrow - x^{2} + 2 x=0\Leftrightarrow x\in\{0;2\}$. $S=\displaystyle\int_{0}^{2}\left|- x^{2} + 2 x\right|\,\mathrm{d}x=\dfrac{4}{3}$.}]{Diện tích hình phẳng giới hạn bởi đồ thị các hàm số $y=- x^{2}$ và $y=- 2 x$ bằng}
  {$\dfrac{4}{3}$}
  {$-\dfrac{4}{3}$}{$\dfrac{4 \pi}{3}$}{$\dfrac{8}{3}$}

\cauTN[id=C4-B13-D02-TH-035,muc=TH,loigiai={$S=\displaystyle\int_{-1}^{1}\left|x^{2} + 1-0\right|\,\mathrm{d}x=\dfrac{8}{3}$.}]{Diện tích hình phẳng giới hạn bởi đồ thị các hàm số $y=x^{2} + 1$, $y=0$ và hai đường thẳng $x=-1$, $x=1$ bằng}
  {$\dfrac{8}{3}$}
  {$\dfrac{8 \pi}{3}$}{$\dfrac{11}{3}$}{$\dfrac{16}{3}$}

\cauTN[id=C4-B13-D02-TH-036,muc=TH,loigiai={$S=\displaystyle\int_{0}^{1}\left|x^{2} + 2-x\right|\,\mathrm{d}x=\dfrac{11}{6}$.}]{Diện tích hình phẳng giới hạn bởi đồ thị các hàm số $y=x^{2} + 2$, $y=x$ và hai đường thẳng $x=0$, $x=1$ bằng}
  {$\dfrac{11}{6}$}
  {$\dfrac{17}{6}$}{$\dfrac{7}{3}$}{$\dfrac{11 \pi}{6}$}

\cauTN[id=C4-B13-D02-TH-037,muc=TH,loigiai={$S=\displaystyle\int_{0}^{\dfrac{\pi}{4}}\left|\sin x-\cos x\right|\,\mathrm{d}x=-1 + \sqrt{2}$.}]{Diện tích hình phẳng giới hạn bởi đồ thị các hàm số $y=\sin x$, $y=\cos x$ và hai đường thẳng $x=0$, $x=\dfrac{\pi}{4}$ bằng}
  {$-1 + \sqrt{2}$}
  {$1$}{$1 - \dfrac{\sqrt{2}}{2}$}{$\pi (-1 + \sqrt{2})$}

\cauTN[id=C4-B13-D02-TH-038,muc=TH,loigiai={Phương trình hoành độ giao điểm: $- 2 x^{2} + 2 x + 2=2 - x^{2}\Leftrightarrow - x^{2} + 2 x=0\Leftrightarrow x\in\{0;2\}$. $S=\displaystyle\int_{0}^{2}\left|- x^{2} + 2 x\right|\,\mathrm{d}x=\dfrac{4}{3}$.}]{Diện tích hình phẳng giới hạn bởi đồ thị các hàm số $y=- 2 x^{2} + 2 x + 2$ và $y=2 - x^{2}$ bằng}
  {$\dfrac{4}{3}$}
  {$-\dfrac{4}{3}$}{$\dfrac{4 \pi}{3}$}{$\dfrac{8}{3}$}

\cauTN[id=C4-B13-D02-TH-039,muc=TH,loigiai={Phương trình hoành độ giao điểm: $- x^{2} - 4 x + 2=2 - x\Leftrightarrow - x^{2} - 3 x=0\Leftrightarrow x\in\{-3;0\}$. $S=\displaystyle\int_{-3}^{0}\left|- x^{2} - 3 x\right|\,\mathrm{d}x=\dfrac{9}{2}$.}]{Diện tích hình phẳng giới hạn bởi đồ thị các hàm số $y=- x^{2} - 4 x + 2$ và $y=2 - x$ bằng}
  {$\dfrac{9}{2}$}
  {$-\dfrac{9}{2}$}{$\dfrac{9 \pi}{2}$}{$9$}

\cauTN[id=C4-B13-D02-TH-040,muc=TH,loigiai={Phương trình hoành độ giao điểm: $- x^{2} - x + 2=x - 1\Leftrightarrow - x^{2} - 2 x + 3=0\Leftrightarrow x\in\{-3;1\}$. $S=\displaystyle\int_{-3}^{1}\left|- x^{2} - 2 x + 3\right|\,\mathrm{d}x=\dfrac{32}{3}$.}]{Diện tích hình phẳng giới hạn bởi đồ thị các hàm số $y=- x^{2} - x + 2$ và $y=x - 1$ bằng}
  {$\dfrac{32}{3}$}
  {$-\dfrac{32}{3}$}{$\dfrac{32 \pi}{3}$}{$\dfrac{64}{3}$}

\cauTN[id=C4-B13-D02-TH-041,muc=TH,loigiai={$S=\displaystyle\int_{0}^{2}\left|x^{2} + 1-0\right|\,\mathrm{d}x=\dfrac{14}{3}$.}]{Diện tích hình phẳng giới hạn bởi đồ thị các hàm số $y=x^{2} + 1$, $y=0$ và hai đường thẳng $x=0$, $x=2$ bằng}
  {$\dfrac{14}{3}$}
  {$\dfrac{14 \pi}{3}$}{$\dfrac{17}{3}$}{$\dfrac{28}{3}$}

\cauTN[id=C4-B13-D02-TH-042,muc=TH,loigiai={$S=\displaystyle\int_{0}^{2}\left|\mathrm{e}^x-x\right|\,\mathrm{d}x=\mathrm{e}^{2} - 3$.}]{Diện tích hình phẳng giới hạn bởi đồ thị các hàm số $y=\mathrm{e}^x$, $y=x$ và hai đường thẳng $x=0$, $x=2$ bằng}
  {$\mathrm{e}^{2} - 3$}
  {$1 + \mathrm{e}^{2}$}{$\mathrm{e}^{2} - 1$}{$\pi (-3 + \mathrm{e}^{2})$}

\cauTN[id=C4-B13-D02-TH-043,muc=TH,loigiai={$S=\displaystyle\int_{0}^{2}\left|x^{2} + 2-x\right|\,\mathrm{d}x=\dfrac{14}{3}$.}]{Diện tích hình phẳng giới hạn bởi đồ thị các hàm số $y=x^{2} + 2$, $y=x$ và hai đường thẳng $x=0$, $x=2$ bằng}
  {$\dfrac{14}{3}$}
  {$\dfrac{26}{3}$}{$\dfrac{20}{3}$}{$\dfrac{14 \pi}{3}$}

\cauTN[id=C4-B13-D02-TH-044,muc=TH,loigiai={$S=\displaystyle\int_{0}^{2}\left|\mathrm{e}^x-(- x)\right|\,\mathrm{d}x=1 + \mathrm{e}^{2}$.}]{Diện tích hình phẳng giới hạn bởi đồ thị các hàm số $y=\mathrm{e}^x$, $y=- x$ và hai đường thẳng $x=0$, $x=2$ bằng}
  {$1 + \mathrm{e}^{2}$}
  {$\mathrm{e}^{2} - 3$}{$\mathrm{e}^{2} - 1$}{$\pi (1 + \mathrm{e}^{2})$}

\cauTN[id=C4-B13-D02-TH-045,muc=TH,loigiai={Phương trình hoành độ giao điểm: $- x^{2} - 3 x - 2=- x - 2\Leftrightarrow - x^{2} - 2 x=0\Leftrightarrow x\in\{-2;0\}$. $S=\displaystyle\int_{-2}^{0}\left|- x^{2} - 2 x\right|\,\mathrm{d}x=\dfrac{4}{3}$.}]{Diện tích hình phẳng giới hạn bởi đồ thị các hàm số $y=- x^{2} - 3 x - 2$ và $y=- x - 2$ bằng}
  {$\dfrac{4}{3}$}
  {$-\dfrac{4}{3}$}{$\dfrac{4 \pi}{3}$}{$\dfrac{8}{3}$}

\cauTN[id=C4-B13-D02-TH-046,muc=TH,loigiai={Phương trình hoành độ giao điểm: $x^{2} + 3 x + 3=x + 3\Leftrightarrow x^{2} + 2 x=0\Leftrightarrow x\in\{-2;0\}$. $S=\displaystyle\int_{-2}^{0}\left|x^{2} + 2 x\right|\,\mathrm{d}x=\dfrac{4}{3}$.}]{Diện tích hình phẳng giới hạn bởi đồ thị các hàm số $y=x^{2} + 3 x + 3$ và $y=x + 3$ bằng}
  {$\dfrac{4}{3}$}
  {$\dfrac{4 \pi}{3}$}{$\dfrac{8}{3}$}{$\dfrac{7}{3}$}

\cauTN[id=C4-B13-D02-TH-047,muc=TH,loigiai={$S=\displaystyle\int_{0}^{2}\left|x^{2} + 2-(- x)\right|\,\mathrm{d}x=\dfrac{26}{3}$.}]{Diện tích hình phẳng giới hạn bởi đồ thị các hàm số $y=x^{2} + 2$, $y=- x$ và hai đường thẳng $x=0$, $x=2$ bằng}
  {$\dfrac{26}{3}$}
  {$\dfrac{14}{3}$}{$\dfrac{20}{3}$}{$\dfrac{26 \pi}{3}$}

\cauTN[id=C4-B13-D02-TH-048,muc=TH,loigiai={Phương trình hoành độ giao điểm: $x^{2} + 4 x - 3=2 x\Leftrightarrow x^{2} + 2 x - 3=0\Leftrightarrow x\in\{-3;1\}$. $S=\displaystyle\int_{-3}^{1}\left|x^{2} + 2 x - 3\right|\,\mathrm{d}x=\dfrac{32}{3}$.}]{Diện tích hình phẳng giới hạn bởi đồ thị các hàm số $y=x^{2} + 4 x - 3$ và $y=2 x$ bằng}
  {$\dfrac{32}{3}$}
  {$\dfrac{32 \pi}{3}$}{$\dfrac{64}{3}$}{$\dfrac{35}{3}$}

\cauDS[id=C4-B13-D02-TH-049,muc=TH,loigiai={$S=-2 + 2 \sqrt{2}$. a) Sai: không có thừa số $\pi$. b) Đúng. c) Sai: đúng là $-2 + 2 \sqrt{2}$. d) Đúng.}]{Gọi $S$ là diện tích hình phẳng giới hạn bởi đồ thị các hàm số $y=\sin x$, $y=\cos x$ và hai đường thẳng $x=0$, $x=\dfrac{\pi}{2}$. Xét tính đúng sai của các mệnh đề sau:}{%
  \y{$S=\pi\displaystyle\int_{0}^{\dfrac{\pi}{2}}\left|\sin x-\cos x\right|\,\mathrm{d}x$.}{s}
  \y{Trên khoảng $(0;\dfrac{\pi}{2})$, hai đồ thị cắt nhau tại điểm có hoành độ $x=\dfrac{\pi}{4}$.}{d}
  \y{$S=2$.}{s}
  \y{$S=-2 + 2 \sqrt{2}$.}{d}}

\cauDS[id=C4-B13-D02-TH-050,muc=TH,loigiai={$S=\dfrac{7}{3}$. a) Sai: đúng là $\dfrac{7}{3}$. b) Đúng. c) Đúng. d) Sai: không có thừa số $\pi$.}]{Gọi $S$ là diện tích hình phẳng giới hạn bởi đồ thị các hàm số $y=x^{2} + 2$, $y=0$ và hai đường thẳng $x=0$, $x=1$. Xét tính đúng sai của các mệnh đề sau:}{%
  \y{$S=\dfrac{7}{3}$.}{s}
  \y{Trên đoạn $[0;1]$, đồ thị $y=x^{2} + 2$ nằm phía trên đồ thị $y=0$.}{d}
  \y{$\displaystyle\int_{0}^{1}x^{2} + 2\,\mathrm{d}x=\dfrac{7}{3}$.}{d}
  \y{$S=\pi\displaystyle\int_{0}^{1}\left|x^{2} + 2-0\right|\,\mathrm{d}x$.}{s}}

\cauDS[id=C4-B13-D02-TH-051,muc=TH,loigiai={$S=2 \sqrt{2}$. a) Đúng. b) Đúng. c) Đúng. d) Sai: không có thừa số $\pi$.}]{Gọi $S$ là diện tích hình phẳng giới hạn bởi đồ thị các hàm số $y=\sin x$, $y=\cos x$ và hai đường thẳng $x=0$, $x=\pi$. Xét tính đúng sai của các mệnh đề sau:}{%
  \y{Trên khoảng $(0;\pi)$, hai đồ thị cắt nhau tại điểm có hoành độ $x=\dfrac{\pi}{4}$.}{d}
  \y{$S=2 \sqrt{2}$.}{d}
  \y{$S=\displaystyle\int_{0}^{\pi}\left|\sin x-\cos x\right|\,\mathrm{d}x$.}{d}
  \y{$S=\pi\displaystyle\int_{0}^{\pi}\left|\sin x-\cos x\right|\,\mathrm{d}x$.}{s}}

\cauDS[id=C4-B13-D02-TH-052,muc=TH,loigiai={$S=\dfrac{14}{3}$. a) Đúng. b) Đúng. c) Sai: đúng là $\dfrac{14}{3}$. d) Sai: không có thừa số $\pi$.}]{Gọi $S$ là diện tích hình phẳng giới hạn bởi đồ thị các hàm số $y=x^{2} + 2$, $y=x$ và hai đường thẳng $x=0$, $x=2$. Xét tính đúng sai của các mệnh đề sau:}{%
  \y{$S=\dfrac{14}{3}$.}{d}
  \y{$S=\displaystyle\int_{0}^{2}\left|x^{2} + 2-x\right|\,\mathrm{d}x$.}{d}
  \y{$S=\dfrac{26}{3}$.}{s}
  \y{$S=\pi\displaystyle\int_{0}^{2}\left|x^{2} + 2-x\right|\,\mathrm{d}x$.}{s}}

\cauDS[id=C4-B13-D02-TH-053,muc=TH,loigiai={$S=- \dfrac{1}{2} + \mathrm{e}$. a) Sai: đúng là $- \dfrac{1}{2} + \mathrm{e}$. b) Đúng. c) Sai: không có thừa số $\pi$. d) Đúng.}]{Gọi $S$ là diện tích hình phẳng giới hạn bởi đồ thị các hàm số $y=\mathrm{e}^x$, $y=- x$ và hai đường thẳng $x=0$, $x=1$. Xét tính đúng sai của các mệnh đề sau:}{%
  \y{$S=- \dfrac{3}{2} + \mathrm{e}$.}{s}
  \y{Trên đoạn $[0;1]$, đồ thị $y=\mathrm{e}^x$ nằm phía trên đồ thị $y=- x$.}{d}
  \y{$S=\pi\displaystyle\int_{0}^{1}\left|\mathrm{e}^x+x\right|\,\mathrm{d}x$.}{s}
  \y{$S=- \dfrac{1}{2} + \mathrm{e}$.}{d}}

\cauDS[id=C4-B13-D02-TH-054,muc=TH,loigiai={$S=2 \sqrt{2}$. a) Đúng. b) Sai: đúng là $2 \sqrt{2}$. c) Đúng. d) Sai: không có thừa số $\pi$.}]{Gọi $S$ là diện tích hình phẳng giới hạn bởi đồ thị các hàm số $y=\sin x$, $y=\cos x$ và hai đường thẳng $x=0$, $x=\pi$. Xét tính đúng sai của các mệnh đề sau:}{%
  \y{$S=\displaystyle\int_{0}^{\pi}\left|\sin x-\cos x\right|\,\mathrm{d}x$.}{d}
  \y{$S=2$.}{s}
  \y{Trên khoảng $(0;\pi)$, hai đồ thị cắt nhau tại điểm có hoành độ $x=\dfrac{\pi}{4}$.}{d}
  \y{$S=\pi\displaystyle\int_{0}^{\pi}\left|\sin x-\cos x\right|\,\mathrm{d}x$.}{s}}

\cauDS[id=C4-B13-D02-TH-055,muc=TH,loigiai={$S=\dfrac{56}{3}$. a) Đúng. b) Đúng. c) Đúng. d) Sai: không có thừa số $\pi$.}]{Gọi $S$ là diện tích hình phẳng giới hạn bởi đồ thị các hàm số $y=x^{2} + 3$, $y=- x$ và hai đường thẳng $x=1$, $x=3$. Xét tính đúng sai của các mệnh đề sau:}{%
  \y{$S=\displaystyle\int_{1}^{3}\left|x^{2} + 3-(- x)\right|\,\mathrm{d}x$.}{d}
  \y{$S=\dfrac{56}{3}$.}{d}
  \y{Trên đoạn $[1;3]$, đồ thị $y=x^{2} + 3$ nằm phía trên đồ thị $y=- x$.}{d}
  \y{$S=\pi\displaystyle\int_{1}^{3}\left|x^{2} + 3+x\right|\,\mathrm{d}x$.}{s}}

\cauDS[id=C4-B13-D02-TH-056,muc=TH,loigiai={$S=-1 + \sqrt{2}$. a) Đúng. b) Đúng. c) Sai: không có thừa số $\pi$. d) Sai: đúng là $-1 + \sqrt{2}$.}]{Gọi $S$ là diện tích hình phẳng giới hạn bởi đồ thị các hàm số $y=\sin x$, $y=\cos x$ và hai đường thẳng $x=0$, $x=\dfrac{\pi}{4}$. Xét tính đúng sai của các mệnh đề sau:}{%
  \y{$\displaystyle\int_{0}^{\dfrac{\pi}{4}}\sin x\,\mathrm{d}x=1 - \dfrac{\sqrt{2}}{2}$.}{d}
  \y{$S=-1 + \sqrt{2}$.}{d}
  \y{$S=\pi\displaystyle\int_{0}^{\dfrac{\pi}{4}}\left|\sin x-\cos x\right|\,\mathrm{d}x$.}{s}
  \y{$S=1$.}{s}}

\cauDS[id=C4-B13-D02-TH-057,muc=TH,loigiai={$S=1 + \sqrt{2}$. a) Đúng. b) Sai: đúng là $1 + \sqrt{2}$. c) Đúng. d) Đúng.}]{Gọi $S$ là diện tích hình phẳng giới hạn bởi đồ thị các hàm số $y=\sin x$, $y=\cos x$ và hai đường thẳng $x=\dfrac{\pi}{4}$, $x=\pi$. Xét tính đúng sai của các mệnh đề sau:}{%
  \y{$\displaystyle\int_{\dfrac{\pi}{4}}^{\pi}\sin x\,\mathrm{d}x=\dfrac{\sqrt{2}}{2} + 1$.}{d}
  \y{$S=1$.}{s}
  \y{$S=1 + \sqrt{2}$.}{d}
  \y{$S=\displaystyle\int_{\dfrac{\pi}{4}}^{\pi}\left|\sin x-\cos x\right|\,\mathrm{d}x$.}{d}}

\cauDS[id=C4-B13-D02-TH-058,muc=TH,loigiai={$S=-1 + \sqrt{2}$. a) Sai: không có thừa số $\pi$. b) Đúng. c) Đúng. d) Đúng.}]{Gọi $S$ là diện tích hình phẳng giới hạn bởi đồ thị các hàm số $y=\sin x$, $y=\cos x$ và hai đường thẳng $x=0$, $x=\dfrac{\pi}{4}$. Xét tính đúng sai của các mệnh đề sau:}{%
  \y{$S=\pi\displaystyle\int_{0}^{\dfrac{\pi}{4}}\left|\sin x-\cos x\right|\,\mathrm{d}x$.}{s}
  \y{Trên đoạn $[0;\dfrac{\pi}{4}]$, đồ thị $y=\sin x$ nằm phía trên đồ thị $y=\cos x$.}{d}
  \y{$\displaystyle\int_{0}^{\dfrac{\pi}{4}}\sin x\,\mathrm{d}x=1 - \dfrac{\sqrt{2}}{2}$.}{d}
  \y{$S=-1 + \sqrt{2}$.}{d}}

\cauDS[id=C4-B13-D02-TH-059,muc=TH,loigiai={$S=2 \sqrt{2}$. a) Đúng. b) Sai: không có thừa số $\pi$. c) Đúng. d) Đúng.}]{Gọi $S$ là diện tích hình phẳng giới hạn bởi đồ thị các hàm số $y=\sin x$, $y=\cos x$ và hai đường thẳng $x=0$, $x=\pi$. Xét tính đúng sai của các mệnh đề sau:}{%
  \y{$S=\displaystyle\int_{0}^{\pi}\left|\sin x-\cos x\right|\,\mathrm{d}x$.}{d}
  \y{$S=\pi\displaystyle\int_{0}^{\pi}\left|\sin x-\cos x\right|\,\mathrm{d}x$.}{s}
  \y{$\displaystyle\int_{0}^{\pi}\sin x\,\mathrm{d}x=2$.}{d}
  \y{Trên khoảng $(0;\pi)$, hai đồ thị cắt nhau tại điểm có hoành độ $x=\dfrac{\pi}{4}$.}{d}}

\cauDS[id=C4-B13-D02-TH-060,muc=TH,loigiai={$S=\dfrac{14}{3}$. a) Đúng. b) Sai: không có thừa số $\pi$. c) Đúng. d) Đúng.}]{Gọi $S$ là diện tích hình phẳng giới hạn bởi đồ thị các hàm số $y=x^{2} + 2$, $y=x$ và hai đường thẳng $x=-1$, $x=1$. Xét tính đúng sai của các mệnh đề sau:}{%
  \y{Trên đoạn $[-1;1]$, đồ thị $y=x^{2} + 2$ nằm phía trên đồ thị $y=x$.}{d}
  \y{$S=\pi\displaystyle\int_{-1}^{1}\left|x^{2} + 2-x\right|\,\mathrm{d}x$.}{s}
  \y{$S=\dfrac{14}{3}$.}{d}
  \y{$\displaystyle\int_{-1}^{1}x^{2} + 2\,\mathrm{d}x=\dfrac{14}{3}$.}{d}}

\cauDS[id=C4-B13-D02-TH-061,muc=TH,loigiai={$S=1 + \sqrt{2}$. a) Đúng. b) Sai: không có thừa số $\pi$. c) Sai: đúng là $1 + \sqrt{2}$. d) Đúng.}]{Gọi $S$ là diện tích hình phẳng giới hạn bởi đồ thị các hàm số $y=\sin x$, $y=\cos x$ và hai đường thẳng $x=\dfrac{\pi}{4}$, $x=\pi$. Xét tính đúng sai của các mệnh đề sau:}{%
  \y{$S=\displaystyle\int_{\dfrac{\pi}{4}}^{\pi}\left|\sin x-\cos x\right|\,\mathrm{d}x$.}{d}
  \y{$S=\pi\displaystyle\int_{\dfrac{\pi}{4}}^{\pi}\left|\sin x-\cos x\right|\,\mathrm{d}x$.}{s}
  \y{$S=1$.}{s}
  \y{$\displaystyle\int_{\dfrac{\pi}{4}}^{\pi}\sin x\,\mathrm{d}x=\dfrac{\sqrt{2}}{2} + 1$.}{d}}

\cauDS[id=C4-B13-D02-TH-062,muc=TH,loigiai={$S=- \mathrm{e} + \dfrac{3}{2} + \mathrm{e}^{2}$. a) Sai: không có thừa số $\pi$. b) Đúng. c) Đúng. d) Đúng.}]{Gọi $S$ là diện tích hình phẳng giới hạn bởi đồ thị các hàm số $y=\mathrm{e}^x$, $y=- x$ và hai đường thẳng $x=1$, $x=2$. Xét tính đúng sai của các mệnh đề sau:}{%
  \y{$S=\pi\displaystyle\int_{1}^{2}\left|\mathrm{e}^x+x\right|\,\mathrm{d}x$.}{s}
  \y{$\displaystyle\int_{1}^{2}\mathrm{e}^x\,\mathrm{d}x=- \mathrm{e} + \mathrm{e}^{2}$.}{d}
  \y{$S=- \mathrm{e} + \dfrac{3}{2} + \mathrm{e}^{2}$.}{d}
  \y{$S=\displaystyle\int_{1}^{2}\left|\mathrm{e}^x-(- x)\right|\,\mathrm{d}x$.}{d}}

\cauDS[id=C4-B13-D02-TH-063,muc=TH,loigiai={$S=-1 + \sqrt{2}$. a) Đúng. b) Đúng. c) Sai: không có thừa số $\pi$. d) Đúng.}]{Gọi $S$ là diện tích hình phẳng giới hạn bởi đồ thị các hàm số $y=\sin x$, $y=\cos x$ và hai đường thẳng $x=0$, $x=\dfrac{\pi}{4}$. Xét tính đúng sai của các mệnh đề sau:}{%
  \y{Trên đoạn $[0;\dfrac{\pi}{4}]$, đồ thị $y=\sin x$ nằm phía trên đồ thị $y=\cos x$.}{d}
  \y{$\displaystyle\int_{0}^{\dfrac{\pi}{4}}\sin x\,\mathrm{d}x=1 - \dfrac{\sqrt{2}}{2}$.}{d}
  \y{$S=\pi\displaystyle\int_{0}^{\dfrac{\pi}{4}}\left|\sin x-\cos x\right|\,\mathrm{d}x$.}{s}
  \y{$S=\displaystyle\int_{0}^{\dfrac{\pi}{4}}\left|\sin x-\cos x\right|\,\mathrm{d}x$.}{d}}

\cauDS[id=C4-B13-D02-TH-064,muc=TH,loigiai={$S=2 \sqrt{2}$. a) Sai: không có thừa số $\pi$. b) Đúng. c) Sai: đúng là $2 \sqrt{2}$. d) Đúng.}]{Gọi $S$ là diện tích hình phẳng giới hạn bởi đồ thị các hàm số $y=\sin x$, $y=\cos x$ và hai đường thẳng $x=0$, $x=\pi$. Xét tính đúng sai của các mệnh đề sau:}{%
  \y{$S=\pi\displaystyle\int_{0}^{\pi}\left|\sin x-\cos x\right|\,\mathrm{d}x$.}{s}
  \y{$\displaystyle\int_{0}^{\pi}\sin x\,\mathrm{d}x=2$.}{d}
  \y{$S=2$.}{s}
  \y{$S=2 \sqrt{2}$.}{d}}

\cauTLN[id=C4-B13-D02-TH-065,muc=TH,loigiai={Phương trình hoành độ giao điểm: $- x^{2} + x + 8=2 x + 2\Leftrightarrow - x^{2} - x + 6=0\Leftrightarrow x\in\{-3;2\}$. $S=\displaystyle\int_{-3}^{2}\left|- x^{2} - x + 6\right|\,\mathrm{d}x=\dfrac{125}{6}$.}]{Tính diện tích hình phẳng giới hạn bởi đồ thị các hàm số $y=- x^{2} + x + 8$ và $y=2 x + 2$ (làm tròn kết quả đến hàng phần mười).}{20,8}

\cauTLN[id=C4-B13-D02-TH-066,muc=TH,loigiai={Phương trình hoành độ giao điểm: $4 - x^{2}=0\Leftrightarrow 4 - x^{2}=0\Leftrightarrow x\in\{-2;2\}$. $S=\displaystyle\int_{-2}^{2}\left|4 - x^{2}\right|\,\mathrm{d}x=\dfrac{32}{3}$.}]{Tính diện tích hình phẳng giới hạn bởi đồ thị các hàm số $y=4 - x^{2}$ và $y=0$ (làm tròn kết quả đến hàng phần mười).}{10,7}

\cauTLN[id=C4-B13-D02-TH-067,muc=TH,loigiai={Phương trình hoành độ giao điểm: $x^{2} + 2 x - 1=2 x\Leftrightarrow x^{2} - 1=0\Leftrightarrow x\in\{-1;1\}$. $S=\displaystyle\int_{-1}^{1}\left|x^{2} - 1\right|\,\mathrm{d}x=\dfrac{4}{3}$.}]{Tính diện tích hình phẳng giới hạn bởi đồ thị các hàm số $y=x^{2} + 2 x - 1$ và $y=2 x$ (làm tròn kết quả đến hàng phần trăm).}{1,33}

\cauTLN[id=C4-B13-D02-TH-068,muc=TH,loigiai={Phương trình hoành độ giao điểm: $- x^{2} + x - 2=-2\Leftrightarrow - x^{2} + x=0\Leftrightarrow x\in\{0;1\}$. $S=\displaystyle\int_{0}^{1}\left|- x^{2} + x\right|\,\mathrm{d}x=\dfrac{1}{6}$.}]{Tính diện tích hình phẳng giới hạn bởi đồ thị các hàm số $y=- x^{2} + x - 2$ và $y=-2$ (làm tròn kết quả đến hàng phần trăm).}{0,17}

\cauTLN[id=C4-B13-D02-TH-069,muc=TH,loigiai={Phương trình hoành độ giao điểm: $- x^{2} - 6 x - 6=- 2 x - 3\Leftrightarrow - x^{2} - 4 x - 3=0\Leftrightarrow x\in\{-3;-1\}$. $S=\displaystyle\int_{-3}^{-1}\left|- x^{2} - 4 x - 3\right|\,\mathrm{d}x=\dfrac{4}{3}$.}]{Tính diện tích hình phẳng giới hạn bởi đồ thị các hàm số $y=- x^{2} - 6 x - 6$ và $y=- 2 x - 3$ (làm tròn kết quả đến hàng phần trăm).}{1,33}

\cauTLN[id=C4-B13-D02-TH-070,muc=TH,loigiai={Phương trình hoành độ giao điểm: $3 x^{2} + 3 x - 3=x^{2} + x + 1\Leftrightarrow 2 x^{2} + 2 x - 4=0\Leftrightarrow x\in\{-2;1\}$. $S=\displaystyle\int_{-2}^{1}\left|2 x^{2} + 2 x - 4\right|\,\mathrm{d}x=9$.}]{Tính diện tích hình phẳng giới hạn bởi đồ thị các hàm số $y=3 x^{2} + 3 x - 3$ và $y=x^{2} + x + 1$.}{9}

\cauTLN[id=C4-B13-D02-TH-071,muc=TH,loigiai={Phương trình hoành độ giao điểm: $- x^{2} - 5 x - 4=2\Leftrightarrow - x^{2} - 5 x - 6=0\Leftrightarrow x\in\{-3;-2\}$. $S=\displaystyle\int_{-3}^{-2}\left|- x^{2} - 5 x - 6\right|\,\mathrm{d}x=\dfrac{1}{6}$.}]{Tính diện tích hình phẳng giới hạn bởi đồ thị các hàm số $y=- x^{2} - 5 x - 4$ và $y=2$ (làm tròn kết quả đến hàng phần trăm).}{0,17}

\cauTLN[id=C4-B13-D02-TH-072,muc=TH,loigiai={Phương trình hoành độ giao điểm: $2 x^{2} - 3 x + 7=x^{2} + 2 x + 1\Leftrightarrow x^{2} - 5 x + 6=0\Leftrightarrow x\in\{2;3\}$. $S=\displaystyle\int_{2}^{3}\left|x^{2} - 5 x + 6\right|\,\mathrm{d}x=\dfrac{1}{6}$.}]{Tính diện tích hình phẳng giới hạn bởi đồ thị các hàm số $y=2 x^{2} - 3 x + 7$ và $y=x^{2} + 2 x + 1$ (làm tròn kết quả đến hàng phần trăm).}{0,17}

\cauTLN[id=C4-B13-D02-TH-073,muc=TH,loigiai={Phương trình hoành độ giao điểm: $x^{2}=2 - x\Leftrightarrow x^{2} + x - 2=0\Leftrightarrow x\in\{-2;1\}$. $S=\displaystyle\int_{-2}^{1}\left|x^{2} + x - 2\right|\,\mathrm{d}x=\dfrac{9}{2}$.}]{Tính diện tích hình phẳng giới hạn bởi đồ thị các hàm số $y=x^{2}$ và $y=2 - x$.}{4,5}

\cauTLN[id=C4-B13-D02-TH-074,muc=TH,loigiai={Phương trình hoành độ giao điểm: $x^{2} - x - 2=0\Leftrightarrow x^{2} - x - 2=0\Leftrightarrow x\in\{-1;2\}$. $S=\displaystyle\int_{-1}^{2}\left|x^{2} - x - 2\right|\,\mathrm{d}x=\dfrac{9}{2}$.}]{Tính diện tích hình phẳng giới hạn bởi đồ thị các hàm số $y=x^{2} - x - 2$ và $y=0$.}{4,5}

\cauTLN[id=C4-B13-D02-TH-075,muc=TH,loigiai={Phương trình hoành độ giao điểm: $- 2 x^{2} - 2 x + 5=- x^{2} - x - 1\Leftrightarrow - x^{2} - x + 6=0\Leftrightarrow x\in\{-3;2\}$. $S=\displaystyle\int_{-3}^{2}\left|- x^{2} - x + 6\right|\,\mathrm{d}x=\dfrac{125}{6}$.}]{Tính diện tích hình phẳng giới hạn bởi đồ thị các hàm số $y=- 2 x^{2} - 2 x + 5$ và $y=- x^{2} - x - 1$ (làm tròn kết quả đến hàng phần mười).}{20,8}

\cauTLN[id=C4-B13-D02-TH-076,muc=TH,loigiai={Phương trình hoành độ giao điểm: $x^{2} + 2 x + 6=3 - 2 x\Leftrightarrow x^{2} + 4 x + 3=0\Leftrightarrow x\in\{-3;-1\}$. $S=\displaystyle\int_{-3}^{-1}\left|x^{2} + 4 x + 3\right|\,\mathrm{d}x=\dfrac{4}{3}$.}]{Tính diện tích hình phẳng giới hạn bởi đồ thị các hàm số $y=x^{2} + 2 x + 6$ và $y=3 - 2 x$ (làm tròn kết quả đến hàng phần trăm).}{1,33}

\cauTLN[id=C4-B13-D02-TH-077,muc=TH,loigiai={Phương trình hoành độ giao điểm: $x^{2} + 2 x - 4=x + 2\Leftrightarrow x^{2} + x - 6=0\Leftrightarrow x\in\{-3;2\}$. $S=\displaystyle\int_{-3}^{2}\left|x^{2} + x - 6\right|\,\mathrm{d}x=\dfrac{125}{6}$.}]{Tính diện tích hình phẳng giới hạn bởi đồ thị các hàm số $y=x^{2} + 2 x - 4$ và $y=x + 2$ (làm tròn kết quả đến hàng phần mười).}{20,8}

\cauTLN[id=C4-B13-D02-TH-078,muc=TH,loigiai={Phương trình hoành độ giao điểm: $2 x^{2} + x - 5=x^{2} + 1\Leftrightarrow x^{2} + x - 6=0\Leftrightarrow x\in\{-3;2\}$. $S=\displaystyle\int_{-3}^{2}\left|x^{2} + x - 6\right|\,\mathrm{d}x=\dfrac{125}{6}$.}]{Tính diện tích hình phẳng giới hạn bởi đồ thị các hàm số $y=2 x^{2} + x - 5$ và $y=x^{2} + 1$ (làm tròn kết quả đến hàng phần mười).}{20,8}

\cauTLN[id=C4-B13-D02-TH-079,muc=TH,loigiai={Phương trình hoành độ giao điểm: $3 x^{2} + 10 x + 14=x^{2} + 2\Leftrightarrow 2 x^{2} + 10 x + 12=0\Leftrightarrow x\in\{-3;-2\}$. $S=\displaystyle\int_{-3}^{-2}\left|2 x^{2} + 10 x + 12\right|\,\mathrm{d}x=\dfrac{1}{3}$.}]{Tính diện tích hình phẳng giới hạn bởi đồ thị các hàm số $y=3 x^{2} + 10 x + 14$ và $y=x^{2} + 2$ (làm tròn kết quả đến hàng phần trăm).}{0,33}

\cauTLN[id=C4-B13-D02-TH-080,muc=TH,loigiai={Phương trình hoành độ giao điểm: $- x^{2} - 3 x + 2=- x - 1\Leftrightarrow - x^{2} - 2 x + 3=0\Leftrightarrow x\in\{-3;1\}$. $S=\displaystyle\int_{-3}^{1}\left|- x^{2} - 2 x + 3\right|\,\mathrm{d}x=\dfrac{32}{3}$.}]{Tính diện tích hình phẳng giới hạn bởi đồ thị các hàm số $y=- x^{2} - 3 x + 2$ và $y=- x - 1$ (làm tròn kết quả đến hàng phần mười).}{10,7}

\cauTN[id=C4-B13-D02-VD-001,muc=VD,loigiai={Phương trình hoành độ giao điểm: $x^{3} - 2 x=- x\Leftrightarrow x^{3} - x=0\Leftrightarrow x\in\{-1;0;1\}$. $S=\displaystyle\int_{-1}^{1}\left|x^{3} - x\right|\,\mathrm{d}x=\dfrac{1}{2}$.}]{Diện tích hình phẳng giới hạn bởi đồ thị các hàm số $y=x^{3} - 2 x$ và $y=- x$ bằng}
  {$\dfrac{1}{2}$}
  {$0$}{$\dfrac{\pi}{2}$}{$1$}

\cauTN[id=C4-B13-D02-VD-002,muc=VD,loigiai={Tiếp tuyến tại $M$: $y=4 x - 4$. $S=\displaystyle\int_{0}^{2}\left[x^2-(4 x - 4)\right]\,\mathrm{d}x=\dfrac{8}{3}$.}]{Diện tích hình phẳng giới hạn bởi parabol $y=x^2$, tiếp tuyến của parabol tại điểm $M(2;4)$ và trục tung bằng}
  {$\dfrac{8}{3}$}
  {$\dfrac{16}{3}$}{$8$}{$\dfrac{11}{3}$}

\cauTN[id=C4-B13-D02-VD-003,muc=VD,loigiai={Phương trình hoành độ giao điểm: $x^{3} - 2 x^{2} - 3 x + 3=x^{2} - 2 x\Leftrightarrow x^{3} - 3 x^{2} - x + 3=0\Leftrightarrow x\in\{-1;1;3\}$. $S=\displaystyle\int_{-1}^{3}\left|x^{3} - 3 x^{2} - x + 3\right|\,\mathrm{d}x=8$.}]{Diện tích hình phẳng giới hạn bởi đồ thị các hàm số $y=x^{3} - 2 x^{2} - 3 x + 3$ và $y=x^{2} - 2 x$ bằng}
  {$8$}
  {$0$}{$8 \pi$}{$16$}

\cauTN[id=C4-B13-D02-VD-004,muc=VD,loigiai={Phương trình hoành độ giao điểm: $- 2 x^{2} + x - 4=- x^{2} - 2 x - 2\Leftrightarrow - x^{2} + 3 x - 2=0\Leftrightarrow x\in\{1;2\}$. $S=\displaystyle\int_{1}^{2}\left|- x^{2} + 3 x - 2\right|\,\mathrm{d}x=\dfrac{1}{6}$.}]{Diện tích hình phẳng giới hạn bởi đồ thị các hàm số $y=- 2 x^{2} + x - 4$ và $y=- x^{2} - 2 x - 2$ bằng}
  {$\dfrac{1}{6}$}
  {$-\dfrac{1}{6}$}{$\dfrac{\pi}{6}$}{$\dfrac{1}{3}$}

\cauTN[id=C4-B13-D02-VD-005,muc=VD,loigiai={Phương trình hoành độ giao điểm: $3 x^{2} - 5 x - 3=x^{2} - x - 3\Leftrightarrow 2 x^{2} - 4 x=0\Leftrightarrow x\in\{0;2\}$. $S=\displaystyle\int_{0}^{2}\left|2 x^{2} - 4 x\right|\,\mathrm{d}x=\dfrac{8}{3}$.}]{Diện tích hình phẳng giới hạn bởi đồ thị các hàm số $y=3 x^{2} - 5 x - 3$ và $y=x^{2} - x - 3$ bằng}
  {$\dfrac{8}{3}$}
  {$-\dfrac{8}{3}$}{$\dfrac{8 \pi}{3}$}{$\dfrac{16}{3}$}

\cauTN[id=C4-B13-D02-VD-006,muc=VD,loigiai={Phương trình hoành độ giao điểm: $x^{3} - 4 x^{2} + 2 x + 6=x\Leftrightarrow x^{3} - 4 x^{2} + x + 6=0\Leftrightarrow x\in\{-1;2;3\}$. $S=\displaystyle\int_{-1}^{3}\left|x^{3} - 4 x^{2} + x + 6\right|\,\mathrm{d}x=\dfrac{71}{6}$.}]{Diện tích hình phẳng giới hạn bởi đồ thị các hàm số $y=x^{3} - 4 x^{2} + 2 x + 6$ và $y=x$ bằng}
  {$\dfrac{71}{6}$}
  {$\dfrac{32}{3}$}{$-\dfrac{32}{3}$}{$\dfrac{71 \pi}{6}$}

\cauTN[id=C4-B13-D02-VD-007,muc=VD,loigiai={Tiếp tuyến tại $M$: $y=6 x - 9$. $S=\displaystyle\int_{0}^{3}\left[x^2-(6 x - 9)\right]\,\mathrm{d}x=9$.}]{Diện tích hình phẳng giới hạn bởi parabol $y=x^2$, tiếp tuyến của parabol tại điểm $M(3;9)$ và trục tung bằng}
  {$9$}
  {$18$}{$27$}{$10$}

\cauTN[id=C4-B13-D02-VD-008,muc=VD,loigiai={Tiếp tuyến tại $M$: $y=2 x - 1$. $S=\displaystyle\int_{0}^{1}\left[x^2-(2 x - 1)\right]\,\mathrm{d}x=\dfrac{1}{3}$.}]{Diện tích hình phẳng giới hạn bởi parabol $y=x^2$, tiếp tuyến của parabol tại điểm $M(1;1)$ và trục tung bằng}
  {$\dfrac{1}{3}$}
  {$\dfrac{2}{3}$}{$1$}{$\dfrac{4}{3}$}

\cauTN[id=C4-B13-D02-VD-009,muc=VD,loigiai={Phương trình hoành độ giao điểm: $x^{3} + 2 x^{2}=x^{2} + 2 x\Leftrightarrow x^{3} + x^{2} - 2 x=0\Leftrightarrow x\in\{-2;0;1\}$. $S=\displaystyle\int_{-2}^{1}\left|x^{3} + x^{2} - 2 x\right|\,\mathrm{d}x=\dfrac{37}{12}$.}]{Diện tích hình phẳng giới hạn bởi đồ thị các hàm số $y=x^{3} + 2 x^{2}$ và $y=x^{2} + 2 x$ bằng}
  {$\dfrac{37}{12}$}
  {$\dfrac{9}{4}$}{$-\dfrac{9}{4}$}{$\dfrac{37 \pi}{12}$}

\cauTN[id=C4-B13-D02-VD-010,muc=VD,loigiai={Phương trình hoành độ giao điểm: $x^{3} - x=0\Leftrightarrow x^{3} - x=0\Leftrightarrow x\in\{-1;0;1\}$. $S=\displaystyle\int_{-1}^{1}\left|x^{3} - x\right|\,\mathrm{d}x=\dfrac{1}{2}$.}]{Diện tích hình phẳng giới hạn bởi đồ thị các hàm số $y=x^{3} - x$ và $y=0$ bằng}
  {$\dfrac{1}{2}$}
  {$0$}{$\dfrac{\pi}{2}$}{$1$}

\cauTN[id=C4-B13-D02-VD-011,muc=VD,loigiai={Phương trình hoành độ giao điểm: $x^{3} - x^{2} - 6 x=- x^{2} - 2 x\Leftrightarrow x^{3} - 4 x=0\Leftrightarrow x\in\{-2;0;2\}$. $S=\displaystyle\int_{-2}^{2}\left|x^{3} - 4 x\right|\,\mathrm{d}x=8$.}]{Diện tích hình phẳng giới hạn bởi đồ thị các hàm số $y=x^{3} - x^{2} - 6 x$ và $y=- x^{2} - 2 x$ bằng}
  {$8$}
  {$0$}{$8 \pi$}{$16$}

\cauTN[id=C4-B13-D02-VD-012,muc=VD,loigiai={Phương trình hoành độ giao điểm: $2 x^{2} - 2 x - 7=x^{2} - x - 1\Leftrightarrow x^{2} - x - 6=0\Leftrightarrow x\in\{-2;3\}$. $S=\displaystyle\int_{-2}^{3}\left|x^{2} - x - 6\right|\,\mathrm{d}x=\dfrac{125}{6}$.}]{Diện tích hình phẳng giới hạn bởi đồ thị các hàm số $y=2 x^{2} - 2 x - 7$ và $y=x^{2} - x - 1$ bằng}
  {$\dfrac{125}{6}$}
  {$-\dfrac{125}{6}$}{$\dfrac{125 \pi}{6}$}{$\dfrac{125}{3}$}

\cauTN[id=C4-B13-D02-VD-013,muc=VD,loigiai={Phương trình hoành độ giao điểm: $x^{3} - 2 x^{2} - x + 3=x^{2}\Leftrightarrow x^{3} - 3 x^{2} - x + 3=0\Leftrightarrow x\in\{-1;1;3\}$. $S=\displaystyle\int_{-1}^{3}\left|x^{3} - 3 x^{2} - x + 3\right|\,\mathrm{d}x=8$.}]{Diện tích hình phẳng giới hạn bởi đồ thị các hàm số $y=x^{3} - 2 x^{2} - x + 3$ và $y=x^{2}$ bằng}
  {$8$}
  {$0$}{$8 \pi$}{$16$}

\cauTN[id=C4-B13-D02-VD-014,muc=VD,loigiai={Phương trình hoành độ giao điểm: $- 2 x^{2} + 4 x - 1=- x^{2} + 2 x - 1\Leftrightarrow - x^{2} + 2 x=0\Leftrightarrow x\in\{0;2\}$. $S=\displaystyle\int_{0}^{2}\left|- x^{2} + 2 x\right|\,\mathrm{d}x=\dfrac{4}{3}$.}]{Diện tích hình phẳng giới hạn bởi đồ thị các hàm số $y=- 2 x^{2} + 4 x - 1$ và $y=- x^{2} + 2 x - 1$ bằng}
  {$\dfrac{4}{3}$}
  {$-\dfrac{4}{3}$}{$\dfrac{4 \pi}{3}$}{$\dfrac{8}{3}$}

\cauTN[id=C4-B13-D02-VD-015,muc=VD,loigiai={Phương trình hoành độ giao điểm: $2 x^{2} + 1=x^{2} - x + 1\Leftrightarrow x^{2} + x=0\Leftrightarrow x\in\{-1;0\}$. $S=\displaystyle\int_{-1}^{0}\left|x^{2} + x\right|\,\mathrm{d}x=\dfrac{1}{6}$.}]{Diện tích hình phẳng giới hạn bởi đồ thị các hàm số $y=2 x^{2} + 1$ và $y=x^{2} - x + 1$ bằng}
  {$\dfrac{1}{6}$}
  {$-\dfrac{1}{6}$}{$\dfrac{\pi}{6}$}{$\dfrac{1}{3}$}

\cauTN[id=C4-B13-D02-VD-016,muc=VD,loigiai={Phương trình hoành độ giao điểm: $x^{3}=x\Leftrightarrow x^{3} - x=0\Leftrightarrow x\in\{-1;0;1\}$. $S=\displaystyle\int_{-1}^{1}\left|x^{3} - x\right|\,\mathrm{d}x=\dfrac{1}{2}$.}]{Diện tích hình phẳng giới hạn bởi đồ thị các hàm số $y=x^{3}$ và $y=x$ bằng}
  {$\dfrac{1}{2}$}
  {$0$}{$\dfrac{\pi}{2}$}{$1$}

\cauTN[id=C4-B13-D02-VD-017,muc=VD,loigiai={Phương trình hoành độ giao điểm: $x^{3} - 3 x=x\Leftrightarrow x^{3} - 4 x=0\Leftrightarrow x\in\{-2;0;2\}$. $S=\displaystyle\int_{-2}^{2}\left|x^{3} - 4 x\right|\,\mathrm{d}x=8$.}]{Diện tích hình phẳng giới hạn bởi đồ thị các hàm số $y=x^{3} - 3 x$ và $y=x$ bằng}
  {$8$}
  {$0$}{$8 \pi$}{$16$}

\cauTN[id=C4-B13-D02-VD-018,muc=VD,loigiai={Phương trình hoành độ giao điểm: $x^{3} - 2 x^{2} - 4 x + 12=x^{2}\Leftrightarrow x^{3} - 3 x^{2} - 4 x + 12=0\Leftrightarrow x\in\{-2;2;3\}$. $S=\displaystyle\int_{-2}^{3}\left|x^{3} - 3 x^{2} - 4 x + 12\right|\,\mathrm{d}x=\dfrac{131}{4}$.}]{Diện tích hình phẳng giới hạn bởi đồ thị các hàm số $y=x^{3} - 2 x^{2} - 4 x + 12$ và $y=x^{2}$ bằng}
  {$\dfrac{131}{4}$}
  {$\dfrac{125}{4}$}{$-\dfrac{125}{4}$}{$\dfrac{131 \pi}{4}$}

\cauTN[id=C4-B13-D02-VD-019,muc=VD,loigiai={Phương trình hoành độ giao điểm: $x^{3} - 4 x=0\Leftrightarrow x^{3} - 4 x=0\Leftrightarrow x\in\{-2;0;2\}$. $S=\displaystyle\int_{-2}^{2}\left|x^{3} - 4 x\right|\,\mathrm{d}x=8$.}]{Diện tích hình phẳng giới hạn bởi đồ thị các hàm số $y=x^{3} - 4 x$ và $y=0$ bằng}
  {$8$}
  {$0$}{$8 \pi$}{$16$}

\cauTN[id=C4-B13-D02-VD-020,muc=VD,loigiai={Phương trình hoành độ giao điểm: $x^{3} - 2 x=- x^{2}\Leftrightarrow x^{3} + x^{2} - 2 x=0\Leftrightarrow x\in\{-2;0;1\}$. $S=\displaystyle\int_{-2}^{1}\left|x^{3} + x^{2} - 2 x\right|\,\mathrm{d}x=\dfrac{37}{12}$.}]{Diện tích hình phẳng giới hạn bởi đồ thị các hàm số $y=x^{3} - 2 x$ và $y=- x^{2}$ bằng}
  {$\dfrac{37}{12}$}
  {$\dfrac{9}{4}$}{$-\dfrac{9}{4}$}{$\dfrac{37 \pi}{12}$}

\cauTN[id=C4-B13-D02-VD-021,muc=VD,loigiai={Phương trình hoành độ giao điểm: $x^{2} + x - 8=- x^{2} + x\Leftrightarrow 2 x^{2} - 8=0\Leftrightarrow x\in\{-2;2\}$. $S=\displaystyle\int_{-2}^{2}\left|2 x^{2} - 8\right|\,\mathrm{d}x=\dfrac{64}{3}$.}]{Diện tích hình phẳng giới hạn bởi đồ thị các hàm số $y=x^{2} + x - 8$ và $y=- x^{2} + x$ bằng}
  {$\dfrac{64}{3}$}
  {$-\dfrac{64}{3}$}{$\dfrac{64 \pi}{3}$}{$\dfrac{128}{3}$}

\cauTN[id=C4-B13-D02-VD-022,muc=VD,loigiai={Phương trình hoành độ giao điểm: $x^{3} - x^{2} - 5 x=x^{2} - 2 x\Leftrightarrow x^{3} - 2 x^{2} - 3 x=0\Leftrightarrow x\in\{-1;0;3\}$. $S=\displaystyle\int_{-1}^{3}\left|x^{3} - 2 x^{2} - 3 x\right|\,\mathrm{d}x=\dfrac{71}{6}$.}]{Diện tích hình phẳng giới hạn bởi đồ thị các hàm số $y=x^{3} - x^{2} - 5 x$ và $y=x^{2} - 2 x$ bằng}
  {$\dfrac{71}{6}$}
  {$-\dfrac{32}{3}$}{$\dfrac{32}{3}$}{$\dfrac{71 \pi}{6}$}

\cauTN[id=C4-B13-D02-VD-023,muc=VD,loigiai={Phương trình hoành độ giao điểm: $x^{3} + x=2 x\Leftrightarrow x^{3} - x=0\Leftrightarrow x\in\{-1;0;1\}$. $S=\displaystyle\int_{-1}^{1}\left|x^{3} - x\right|\,\mathrm{d}x=\dfrac{1}{2}$.}]{Diện tích hình phẳng giới hạn bởi đồ thị các hàm số $y=x^{3} + x$ và $y=2 x$ bằng}
  {$\dfrac{1}{2}$}
  {$0$}{$\dfrac{\pi}{2}$}{$1$}

\cauTN[id=C4-B13-D02-VD-024,muc=VD,loigiai={Phương trình hoành độ giao điểm: $x^{3} - x^{2} - 4 x=x^{2} - x\Leftrightarrow x^{3} - 2 x^{2} - 3 x=0\Leftrightarrow x\in\{-1;0;3\}$. $S=\displaystyle\int_{-1}^{3}\left|x^{3} - 2 x^{2} - 3 x\right|\,\mathrm{d}x=\dfrac{71}{6}$.}]{Diện tích hình phẳng giới hạn bởi đồ thị các hàm số $y=x^{3} - x^{2} - 4 x$ và $y=x^{2} - x$ bằng}
  {$\dfrac{71}{6}$}
  {$-\dfrac{32}{3}$}{$\dfrac{32}{3}$}{$\dfrac{71 \pi}{6}$}

\cauTN[id=C4-B13-D02-VD-025,muc=VD,loigiai={Phương trình hoành độ giao điểm: $2 x^{2} + 2 x - 7=x^{2} + 2 x + 2\Leftrightarrow x^{2} - 9=0\Leftrightarrow x\in\{-3;3\}$. $S=\displaystyle\int_{-3}^{3}\left|x^{2} - 9\right|\,\mathrm{d}x=36$.}]{Diện tích hình phẳng giới hạn bởi đồ thị các hàm số $y=2 x^{2} + 2 x - 7$ và $y=x^{2} + 2 x + 2$ bằng}
  {$36$}
  {$-36$}{$36 \pi$}{$72$}

\cauTN[id=C4-B13-D02-VD-026,muc=VD,loigiai={Phương trình hoành độ giao điểm: $x^{3} + x^{2} - 3 x=- x\Leftrightarrow x^{3} + x^{2} - 2 x=0\Leftrightarrow x\in\{-2;0;1\}$. $S=\displaystyle\int_{-2}^{1}\left|x^{3} + x^{2} - 2 x\right|\,\mathrm{d}x=\dfrac{37}{12}$.}]{Diện tích hình phẳng giới hạn bởi đồ thị các hàm số $y=x^{3} + x^{2} - 3 x$ và $y=- x$ bằng}
  {$\dfrac{37}{12}$}
  {$\dfrac{9}{4}$}{$-\dfrac{9}{4}$}{$\dfrac{37 \pi}{12}$}

\cauTN[id=C4-B13-D02-VD-027,muc=VD,loigiai={Phương trình hoành độ giao điểm: $x^{3} - 3 x^{2} + 5 x=x^{2} + 2 x\Leftrightarrow x^{3} - 4 x^{2} + 3 x=0\Leftrightarrow x\in\{0;1;3\}$. $S=\displaystyle\int_{0}^{3}\left|x^{3} - 4 x^{2} + 3 x\right|\,\mathrm{d}x=\dfrac{37}{12}$.}]{Diện tích hình phẳng giới hạn bởi đồ thị các hàm số $y=x^{3} - 3 x^{2} + 5 x$ và $y=x^{2} + 2 x$ bằng}
  {$\dfrac{37}{12}$}
  {$-\dfrac{9}{4}$}{$\dfrac{9}{4}$}{$\dfrac{37 \pi}{12}$}

\cauTN[id=C4-B13-D02-VD-028,muc=VD,loigiai={Phương trình hoành độ giao điểm: $x^{3} - 2 x=2 x\Leftrightarrow x^{3} - 4 x=0\Leftrightarrow x\in\{-2;0;2\}$. $S=\displaystyle\int_{-2}^{2}\left|x^{3} - 4 x\right|\,\mathrm{d}x=8$.}]{Diện tích hình phẳng giới hạn bởi đồ thị các hàm số $y=x^{3} - 2 x$ và $y=2 x$ bằng}
  {$8$}
  {$0$}{$8 \pi$}{$16$}

\cauTN[id=C4-B13-D02-VD-029,muc=VD,loigiai={Phương trình hoành độ giao điểm: $x^{3} - 5 x=- x\Leftrightarrow x^{3} - 4 x=0\Leftrightarrow x\in\{-2;0;2\}$. $S=\displaystyle\int_{-2}^{2}\left|x^{3} - 4 x\right|\,\mathrm{d}x=8$.}]{Diện tích hình phẳng giới hạn bởi đồ thị các hàm số $y=x^{3} - 5 x$ và $y=- x$ bằng}
  {$8$}
  {$0$}{$8 \pi$}{$16$}

\cauTN[id=C4-B13-D02-VD-030,muc=VD,loigiai={Phương trình hoành độ giao điểm: $x^{3} - 4 x^{2} + 6=- x\Leftrightarrow x^{3} - 4 x^{2} + x + 6=0\Leftrightarrow x\in\{-1;2;3\}$. $S=\displaystyle\int_{-1}^{3}\left|x^{3} - 4 x^{2} + x + 6\right|\,\mathrm{d}x=\dfrac{71}{6}$.}]{Diện tích hình phẳng giới hạn bởi đồ thị các hàm số $y=x^{3} - 4 x^{2} + 6$ và $y=- x$ bằng}
  {$\dfrac{71}{6}$}
  {$\dfrac{32}{3}$}{$-\dfrac{32}{3}$}{$\dfrac{71 \pi}{6}$}

\cauTN[id=C4-B13-D02-VD-031,muc=VD,loigiai={Phương trình hoành độ giao điểm: $x^{3} - 5 x^{2} + 11 x - 6=x^{2}\Leftrightarrow x^{3} - 6 x^{2} + 11 x - 6=0\Leftrightarrow x\in\{1;2;3\}$. $S=\displaystyle\int_{1}^{3}\left|x^{3} - 6 x^{2} + 11 x - 6\right|\,\mathrm{d}x=\dfrac{1}{2}$.}]{Diện tích hình phẳng giới hạn bởi đồ thị các hàm số $y=x^{3} - 5 x^{2} + 11 x - 6$ và $y=x^{2}$ bằng}
  {$\dfrac{1}{2}$}
  {$0$}{$\dfrac{\pi}{2}$}{$1$}

\cauTN[id=C4-B13-D02-VD-032,muc=VD,loigiai={Phương trình hoành độ giao điểm: $x^{3} + 2 x^{2} - 5 x - 4=x^{2} - x\Leftrightarrow x^{3} + x^{2} - 4 x - 4=0\Leftrightarrow x\in\{-2;-1;2\}$. $S=\displaystyle\int_{-2}^{2}\left|x^{3} + x^{2} - 4 x - 4\right|\,\mathrm{d}x=\dfrac{71}{6}$.}]{Diện tích hình phẳng giới hạn bởi đồ thị các hàm số $y=x^{3} + 2 x^{2} - 5 x - 4$ và $y=x^{2} - x$ bằng}
  {$\dfrac{71}{6}$}
  {$-\dfrac{32}{3}$}{$\dfrac{32}{3}$}{$\dfrac{71 \pi}{6}$}

\cauTN[id=C4-B13-D02-VD-033,muc=VD,loigiai={Phương trình hoành độ giao điểm: $2 x^{2} + 5 x + 3=x^{2} + 2 x + 1\Leftrightarrow x^{2} + 3 x + 2=0\Leftrightarrow x\in\{-2;-1\}$. $S=\displaystyle\int_{-2}^{-1}\left|x^{2} + 3 x + 2\right|\,\mathrm{d}x=\dfrac{1}{6}$.}]{Diện tích hình phẳng giới hạn bởi đồ thị các hàm số $y=2 x^{2} + 5 x + 3$ và $y=x^{2} + 2 x + 1$ bằng}
  {$\dfrac{1}{6}$}
  {$-\dfrac{1}{6}$}{$\dfrac{\pi}{6}$}{$\dfrac{1}{3}$}

\cauTN[id=C4-B13-D02-VD-034,muc=VD,loigiai={Phương trình hoành độ giao điểm: $x^{2} - 3 x - 14=- x^{2} - x - 2\Leftrightarrow 2 x^{2} - 2 x - 12=0\Leftrightarrow x\in\{-2;3\}$. $S=\displaystyle\int_{-2}^{3}\left|2 x^{2} - 2 x - 12\right|\,\mathrm{d}x=\dfrac{125}{3}$.}]{Diện tích hình phẳng giới hạn bởi đồ thị các hàm số $y=x^{2} - 3 x - 14$ và $y=- x^{2} - x - 2$ bằng}
  {$\dfrac{125}{3}$}
  {$-\dfrac{125}{3}$}{$\dfrac{125 \pi}{3}$}{$\dfrac{250}{3}$}

\cauTN[id=C4-B13-D02-VD-035,muc=VD,loigiai={Phương trình hoành độ giao điểm: $- 2 x^{2} + 2 x + 8=- x^{2} + 2 x - 1\Leftrightarrow 9 - x^{2}=0\Leftrightarrow x\in\{-3;3\}$. $S=\displaystyle\int_{-3}^{3}\left|9 - x^{2}\right|\,\mathrm{d}x=36$.}]{Diện tích hình phẳng giới hạn bởi đồ thị các hàm số $y=- 2 x^{2} + 2 x + 8$ và $y=- x^{2} + 2 x - 1$ bằng}
  {$36$}
  {$-36$}{$36 \pi$}{$72$}

\cauTN[id=C4-B13-D02-VD-036,muc=VD,loigiai={Phương trình hoành độ giao điểm: $x^{3} - x^{2} - 2 x=- x^{2} - x\Leftrightarrow x^{3} - x=0\Leftrightarrow x\in\{-1;0;1\}$. $S=\displaystyle\int_{-1}^{1}\left|x^{3} - x\right|\,\mathrm{d}x=\dfrac{1}{2}$.}]{Diện tích hình phẳng giới hạn bởi đồ thị các hàm số $y=x^{3} - x^{2} - 2 x$ và $y=- x^{2} - x$ bằng}
  {$\dfrac{1}{2}$}
  {$0$}{$\dfrac{\pi}{2}$}{$1$}

\cauDS[id=C4-B13-D02-VD-037,muc=VD,loigiai={Phương trình hoành độ giao điểm: $x^{3} - 4 x^{2} - 6 x + 12=- x^{2} - 2 x\Leftrightarrow x^{3} - 3 x^{2} - 4 x + 12=0\Leftrightarrow x\in\{-2;2;3\}$. $S=\displaystyle\int_{-2}^{3}\left|x^{3} - 3 x^{2} - 4 x + 12\right|\,\mathrm{d}x=\dfrac{131}{4}$. a) Đúng. b) Sai: có $3$ điểm chung. c) Sai: đúng là $\dfrac{131}{4}$. d) Sai: đúng là $\dfrac{131}{4}$.}]{Gọi $S$ là diện tích hình phẳng giới hạn bởi đồ thị các hàm số $y=x^{3} - 4 x^{2} - 6 x + 12$ và $y=- x^{2} - 2 x$. Xét tính đúng sai của các mệnh đề sau:}{%
  \y{$S=\displaystyle\int_{-2}^{3}\left|x^{3} - 3 x^{2} - 4 x + 12\right|\,\mathrm{d}x$.}{d}
  \y{Hai đồ thị có 4 điểm chung.}{s}
  \y{$S=\dfrac{139}{4}$.}{s}
  \y{$S=\dfrac{125}{4}$.}{s}}

\cauDS[id=C4-B13-D02-VD-038,muc=VD,loigiai={Phương trình hoành độ giao điểm: $x^{3} - 5 x^{2} + 7 x=x\Leftrightarrow x^{3} - 5 x^{2} + 6 x=0\Leftrightarrow x\in\{0;2;3\}$. $S=\displaystyle\int_{0}^{3}\left|x^{3} - 5 x^{2} + 6 x\right|\,\mathrm{d}x=\dfrac{37}{12}$. a) Đúng. b) Sai: đúng là $\dfrac{37}{12}$. c) Sai: đúng là $\dfrac{37}{12}$. d) Đúng.}]{Gọi $S$ là diện tích hình phẳng giới hạn bởi đồ thị các hàm số $y=x^{3} - 5 x^{2} + 7 x$ và $y=x$. Xét tính đúng sai của các mệnh đề sau:}{%
  \y{$S=\displaystyle\int_{0}^{3}\left|x^{3} - 5 x^{2} + 6 x\right|\,\mathrm{d}x$.}{d}
  \y{$S=\dfrac{61}{12}$.}{s}
  \y{$S=\dfrac{9}{4}$.}{s}
  \y{Hai đồ thị có 3 điểm chung.}{d}}

\cauDS[id=C4-B13-D02-VD-039,muc=VD,loigiai={Phương trình hoành độ giao điểm: $x^{2} + 4 x + 7=- x^{2} - 2 x + 3\Leftrightarrow 2 x^{2} + 6 x + 4=0\Leftrightarrow x\in\{-2;-1\}$. $S=\displaystyle\int_{-2}^{-1}\left|2 x^{2} + 6 x + 4\right|\,\mathrm{d}x=\dfrac{1}{3}$. a) Đúng. b) Sai: đúng là $\dfrac{1}{3}$. c) Đúng: đúng là $\dfrac{1}{3}$. d) Sai: có $2$ điểm chung.}]{Gọi $S$ là diện tích hình phẳng giới hạn bởi đồ thị các hàm số $y=x^{2} + 4 x + 7$ và $y=- x^{2} - 2 x + 3$. Xét tính đúng sai của các mệnh đề sau:}{%
  \y{$S=\displaystyle\int_{-2}^{-1}\left|2 x^{2} + 6 x + 4\right|\,\mathrm{d}x$.}{d}
  \y{$S=\dfrac{7}{3}$.}{s}
  \y{$S=\dfrac{1}{3}$.}{d}
  \y{Hai đồ thị có 3 điểm chung.}{s}}

\cauDS[id=C4-B13-D02-VD-040,muc=VD,loigiai={Phương trình hoành độ giao điểm: $- 2 x^{2} - 3 x + 8=- x^{2} - 2 x + 2\Leftrightarrow - x^{2} - x + 6=0\Leftrightarrow x\in\{-3;2\}$. $S=\displaystyle\int_{-3}^{2}\left|- x^{2} - x + 6\right|\,\mathrm{d}x=\dfrac{125}{6}$. a) Đúng. b) Sai: đúng là $\dfrac{125}{6}$. c) Sai: có $2$ điểm chung. d) Đúng.}]{Gọi $S$ là diện tích hình phẳng giới hạn bởi đồ thị các hàm số $y=- 2 x^{2} - 3 x + 8$ và $y=- x^{2} - 2 x + 2$. Xét tính đúng sai của các mệnh đề sau:}{%
  \y{Hai đồ thị có 2 điểm chung.}{d}
  \y{$S=\dfrac{137}{6}$.}{s}
  \y{Hai đồ thị có 3 điểm chung.}{s}
  \y{$S=\dfrac{125}{6}$.}{d}}

\cauDS[id=C4-B13-D02-VD-041,muc=VD,loigiai={Phương trình hoành độ giao điểm: $x^{3} - 2 x^{2} - 2 x + 2=- x\Leftrightarrow x^{3} - 2 x^{2} - x + 2=0\Leftrightarrow x\in\{-1;1;2\}$. $S=\displaystyle\int_{-1}^{2}\left|x^{3} - 2 x^{2} - x + 2\right|\,\mathrm{d}x=\dfrac{37}{12}$. a) Sai: đúng là $\dfrac{37}{12}$. b) Đúng. c) Sai: có $3$ điểm chung. d) Sai: đúng là $\dfrac{37}{12}$.}]{Gọi $S$ là diện tích hình phẳng giới hạn bởi đồ thị các hàm số $y=x^{3} - 2 x^{2} - 2 x + 2$ và $y=- x$. Xét tính đúng sai của các mệnh đề sau:}{%
  \y{$S=\dfrac{9}{4}$.}{s}
  \y{Hai đồ thị có 3 điểm chung.}{d}
  \y{Hai đồ thị có 4 điểm chung.}{s}
  \y{$S=\dfrac{61}{12}$.}{s}}

\cauDS[id=C4-B13-D02-VD-042,muc=VD,loigiai={Phương trình hoành độ giao điểm: $x^{3} - 3 x^{2} + 2 x + 6=x^{2} + x\Leftrightarrow x^{3} - 4 x^{2} + x + 6=0\Leftrightarrow x\in\{-1;2;3\}$. $S=\displaystyle\int_{-1}^{3}\left|x^{3} - 4 x^{2} + x + 6\right|\,\mathrm{d}x=\dfrac{71}{6}$. a) Sai: có $3$ điểm chung. b) Đúng. c) Sai: đúng là $\dfrac{71}{6}$. d) Đúng.}]{Gọi $S$ là diện tích hình phẳng giới hạn bởi đồ thị các hàm số $y=x^{3} - 3 x^{2} + 2 x + 6$ và $y=x^{2} + x$. Xét tính đúng sai của các mệnh đề sau:}{%
  \y{Hai đồ thị có 4 điểm chung.}{s}
  \y{$S=\dfrac{71}{6}$.}{d}
  \y{$S=\dfrac{32}{3}$.}{s}
  \y{$S=\displaystyle\int_{-1}^{3}\left|x^{3} - 4 x^{2} + x + 6\right|\,\mathrm{d}x$.}{d}}

\cauDS[id=C4-B13-D02-VD-043,muc=VD,loigiai={Phương trình hoành độ giao điểm: $x^{3} - 2 x=- x\Leftrightarrow x^{3} - x=0\Leftrightarrow x\in\{-1;0;1\}$. $S=\displaystyle\int_{-1}^{1}\left|x^{3} - x\right|\,\mathrm{d}x=\dfrac{1}{2}$. a) Đúng. b) Sai: đúng là $\dfrac{1}{2}$. c) Sai: đúng là $\dfrac{1}{2}$. d) Đúng.}]{Gọi $S$ là diện tích hình phẳng giới hạn bởi đồ thị các hàm số $y=x^{3} - 2 x$ và $y=- x$. Xét tính đúng sai của các mệnh đề sau:}{%
  \y{$S=\displaystyle\int_{-1}^{1}\left|x^{3} - x\right|\,\mathrm{d}x$.}{d}
  \y{$S=0$.}{s}
  \y{$S=\dfrac{5}{2}$.}{s}
  \y{$S=\dfrac{1}{2}$.}{d}}

\cauDS[id=C4-B13-D02-VD-044,muc=VD,loigiai={Phương trình hoành độ giao điểm: $x^{3} - 4 x^{2} + x + 3=- x^{2} + 2 x\Leftrightarrow x^{3} - 3 x^{2} - x + 3=0\Leftrightarrow x\in\{-1;1;3\}$. $S=\displaystyle\int_{-1}^{3}\left|x^{3} - 3 x^{2} - x + 3\right|\,\mathrm{d}x=8$. a) Đúng. b) Sai: đúng là $8$. c) Đúng. d) Sai: có $3$ điểm chung.}]{Gọi $S$ là diện tích hình phẳng giới hạn bởi đồ thị các hàm số $y=x^{3} - 4 x^{2} + x + 3$ và $y=- x^{2} + 2 x$. Xét tính đúng sai của các mệnh đề sau:}{%
  \y{$S=\displaystyle\int_{-1}^{3}\left|x^{3} - 3 x^{2} - x + 3\right|\,\mathrm{d}x$.}{d}
  \y{$S=10$.}{s}
  \y{Hai đồ thị có 3 điểm chung.}{d}
  \y{Hai đồ thị có 4 điểm chung.}{s}}

\cauDS[id=C4-B13-D02-VD-045,muc=VD,loigiai={Phương trình hoành độ giao điểm: $x^{3} - 2 x=2 x\Leftrightarrow x^{3} - 4 x=0\Leftrightarrow x\in\{-2;0;2\}$. $S=\displaystyle\int_{-2}^{2}\left|x^{3} - 4 x\right|\,\mathrm{d}x=8$. a) Đúng. b) Đúng. c) Đúng. d) Sai: đúng là $8$.}]{Gọi $S$ là diện tích hình phẳng giới hạn bởi đồ thị các hàm số $y=x^{3} - 2 x$ và $y=2 x$. Xét tính đúng sai của các mệnh đề sau:}{%
  \y{$S=\displaystyle\int_{-2}^{2}\left|x^{3} - 4 x\right|\,\mathrm{d}x$.}{d}
  \y{Hai đồ thị có 3 điểm chung.}{d}
  \y{$S=8$.}{d}
  \y{$S=10$.}{s}}

\cauDS[id=C4-B13-D02-VD-046,muc=VD,loigiai={Phương trình hoành độ giao điểm: $x^{3} + 2 x^{2} + x - 2=2 x\Leftrightarrow x^{3} + 2 x^{2} - x - 2=0\Leftrightarrow x\in\{-2;-1;1\}$. $S=\displaystyle\int_{-2}^{1}\left|x^{3} + 2 x^{2} - x - 2\right|\,\mathrm{d}x=\dfrac{37}{12}$. a) Sai: có $3$ điểm chung. b) Sai: đúng là $\dfrac{37}{12}$. c) Sai: đúng là $\dfrac{37}{12}$. d) Đúng.}]{Gọi $S$ là diện tích hình phẳng giới hạn bởi đồ thị các hàm số $y=x^{3} + 2 x^{2} + x - 2$ và $y=2 x$. Xét tính đúng sai của các mệnh đề sau:}{%
  \y{Hai đồ thị có 4 điểm chung.}{s}
  \y{$S=\dfrac{9}{4}$.}{s}
  \y{$S=\dfrac{61}{12}$.}{s}
  \y{$S=\displaystyle\int_{-2}^{1}\left|x^{3} + 2 x^{2} - x - 2\right|\,\mathrm{d}x$.}{d}}

\cauDS[id=C4-B13-D02-VD-047,muc=VD,loigiai={Phương trình hoành độ giao điểm: $2 x^{2} + 3 x + 8=x^{2} - 2 x + 2\Leftrightarrow x^{2} + 5 x + 6=0\Leftrightarrow x\in\{-3;-2\}$. $S=\displaystyle\int_{-3}^{-2}\left|x^{2} + 5 x + 6\right|\,\mathrm{d}x=\dfrac{1}{6}$. a) Đúng. b) Đúng. c) Đúng: đúng là $\dfrac{1}{6}$. d) Sai: đúng là $\dfrac{1}{6}$.}]{Gọi $S$ là diện tích hình phẳng giới hạn bởi đồ thị các hàm số $y=2 x^{2} + 3 x + 8$ và $y=x^{2} - 2 x + 2$. Xét tính đúng sai của các mệnh đề sau:}{%
  \y{$S=\displaystyle\int_{-3}^{-2}\left|x^{2} + 5 x + 6\right|\,\mathrm{d}x$.}{d}
  \y{Hai đồ thị có 2 điểm chung.}{d}
  \y{$S=\dfrac{1}{6}$.}{d}
  \y{$S=\dfrac{13}{6}$.}{s}}

\cauDS[id=C4-B13-D02-VD-048,muc=VD,loigiai={Phương trình hoành độ giao điểm: $x^{3} - 3 x^{2} - 3 x + 3=- 2 x\Leftrightarrow x^{3} - 3 x^{2} - x + 3=0\Leftrightarrow x\in\{-1;1;3\}$. $S=\displaystyle\int_{-1}^{3}\left|x^{3} - 3 x^{2} - x + 3\right|\,\mathrm{d}x=8$. a) Đúng. b) Đúng. c) Sai: có $3$ điểm chung. d) Sai: đúng là $8$.}]{Gọi $S$ là diện tích hình phẳng giới hạn bởi đồ thị các hàm số $y=x^{3} - 3 x^{2} - 3 x + 3$ và $y=- 2 x$. Xét tính đúng sai của các mệnh đề sau:}{%
  \y{$S=8$.}{d}
  \y{Hai đồ thị có 3 điểm chung.}{d}
  \y{Hai đồ thị có 4 điểm chung.}{s}
  \y{$S=0$.}{s}}

\cauTLN[id=C4-B13-D02-VD-049,muc=VD,loigiai={Phương trình hoành độ giao điểm: $x^{3} - 3 x=x\Leftrightarrow x^{3} - 4 x=0\Leftrightarrow x\in\{-2;0;2\}$. $S=\displaystyle\int_{-2}^{2}\left|x^{3} - 4 x\right|\,\mathrm{d}x=8$.}]{Tính diện tích hình phẳng giới hạn bởi đồ thị các hàm số $y=x^{3} - 3 x$ và $y=x$.}{8}

\cauTLN[id=C4-B13-D02-VD-050,muc=VD,loigiai={Phương trình hoành độ giao điểm: $x^{3}=x\Leftrightarrow x^{3} - x=0\Leftrightarrow x\in\{-1;0;1\}$. $S=\displaystyle\int_{-1}^{1}\left|x^{3} - x\right|\,\mathrm{d}x=\dfrac{1}{2}$.}]{Tính diện tích hình phẳng giới hạn bởi đồ thị các hàm số $y=x^{3}$ và $y=x$.}{0,5}

\cauTLN[id=C4-B13-D02-VD-051,muc=VD,loigiai={Phương trình hoành độ giao điểm: $x^{3} - 5 x=- x\Leftrightarrow x^{3} - 4 x=0\Leftrightarrow x\in\{-2;0;2\}$. $S=\displaystyle\int_{-2}^{2}\left|x^{3} - 4 x\right|\,\mathrm{d}x=8$.}]{Tính diện tích hình phẳng giới hạn bởi đồ thị các hàm số $y=x^{3} - 5 x$ và $y=- x$.}{8}

\cauTLN[id=C4-B13-D02-VD-052,muc=VD,loigiai={Phương trình hoành độ giao điểm: $x^{3} - x=- x^{2} + x\Leftrightarrow x^{3} + x^{2} - 2 x=0\Leftrightarrow x\in\{-2;0;1\}$. $S=\displaystyle\int_{-2}^{1}\left|x^{3} + x^{2} - 2 x\right|\,\mathrm{d}x=\dfrac{37}{12}$.}]{Tính diện tích hình phẳng giới hạn bởi đồ thị các hàm số $y=x^{3} - x$ và $y=- x^{2} + x$ (làm tròn kết quả đến hàng phần trăm).}{3,08}

\cauTLN[id=C4-B13-D02-VD-053,muc=VD,loigiai={Phương trình hoành độ giao điểm: $x^{3} - 2 x=- x\Leftrightarrow x^{3} - x=0\Leftrightarrow x\in\{-1;0;1\}$. $S=\displaystyle\int_{-1}^{1}\left|x^{3} - x\right|\,\mathrm{d}x=\dfrac{1}{2}$.}]{Tính diện tích hình phẳng giới hạn bởi đồ thị các hàm số $y=x^{3} - 2 x$ và $y=- x$.}{0,5}

\cauTLN[id=C4-B13-D02-VD-054,muc=VD,loigiai={Phương trình hoành độ giao điểm: $3 x^{2} - 12=x^{2} + 2 x\Leftrightarrow 2 x^{2} - 2 x - 12=0\Leftrightarrow x\in\{-2;3\}$. $S=\displaystyle\int_{-2}^{3}\left|2 x^{2} - 2 x - 12\right|\,\mathrm{d}x=\dfrac{125}{3}$.}]{Tính diện tích hình phẳng giới hạn bởi đồ thị các hàm số $y=3 x^{2} - 12$ và $y=x^{2} + 2 x$ (làm tròn kết quả đến hàng phần mười).}{41,7}

\cauTLN[id=C4-B13-D02-VD-055,muc=VD,loigiai={Phương trình hoành độ giao điểm: $x^{3} - x^{2} - 4 x + 4=0\Leftrightarrow x^{3} - x^{2} - 4 x + 4=0\Leftrightarrow x\in\{-2;1;2\}$. $S=\displaystyle\int_{-2}^{2}\left|x^{3} - x^{2} - 4 x + 4\right|\,\mathrm{d}x=\dfrac{71}{6}$.}]{Tính diện tích hình phẳng giới hạn bởi đồ thị các hàm số $y=x^{3} - x^{2} - 4 x + 4$ và $y=0$ (làm tròn kết quả đến hàng phần mười).}{11,8}

\cauTLN[id=C4-B13-D02-VD-056,muc=VD,loigiai={Phương trình hoành độ giao điểm: $3 x^{2} - 10=x^{2} - 2\Leftrightarrow 2 x^{2} - 8=0\Leftrightarrow x\in\{-2;2\}$. $S=\displaystyle\int_{-2}^{2}\left|2 x^{2} - 8\right|\,\mathrm{d}x=\dfrac{64}{3}$.}]{Tính diện tích hình phẳng giới hạn bởi đồ thị các hàm số $y=3 x^{2} - 10$ và $y=x^{2} - 2$ (làm tròn kết quả đến hàng phần mười).}{21,3}

\cauTLN[id=C4-B13-D02-VD-057,muc=VD,loigiai={Phương trình hoành độ giao điểm: $3 x^{2} - 3 x + 1=x^{2} + x + 1\Leftrightarrow 2 x^{2} - 4 x=0\Leftrightarrow x\in\{0;2\}$. $S=\displaystyle\int_{0}^{2}\left|2 x^{2} - 4 x\right|\,\mathrm{d}x=\dfrac{8}{3}$.}]{Tính diện tích hình phẳng giới hạn bởi đồ thị các hàm số $y=3 x^{2} - 3 x + 1$ và $y=x^{2} + x + 1$ (làm tròn kết quả đến hàng phần trăm).}{2,67}

\cauTLN[id=C4-B13-D02-VD-058,muc=VD,loigiai={Phương trình hoành độ giao điểm: $2 x^{2} + 5 x=x^{2} + 2 x\Leftrightarrow x^{2} + 3 x=0\Leftrightarrow x\in\{-3;0\}$. $S=\displaystyle\int_{-3}^{0}\left|x^{2} + 3 x\right|\,\mathrm{d}x=\dfrac{9}{2}$.}]{Tính diện tích hình phẳng giới hạn bởi đồ thị các hàm số $y=2 x^{2} + 5 x$ và $y=x^{2} + 2 x$.}{4,5}

\cauTLN[id=C4-B13-D02-VD-059,muc=VD,loigiai={Phương trình hoành độ giao điểm: $2 x^{2} - 3 x + 2=x^{2}\Leftrightarrow x^{2} - 3 x + 2=0\Leftrightarrow x\in\{1;2\}$. $S=\displaystyle\int_{1}^{2}\left|x^{2} - 3 x + 2\right|\,\mathrm{d}x=\dfrac{1}{6}$.}]{Tính diện tích hình phẳng giới hạn bởi đồ thị các hàm số $y=2 x^{2} - 3 x + 2$ và $y=x^{2}$ (làm tròn kết quả đến hàng phần trăm).}{0,17}

\cauTLN[id=C4-B13-D02-VD-060,muc=VD,loigiai={Phương trình hoành độ giao điểm: $x^{3} - x=0\Leftrightarrow x^{3} - x=0\Leftrightarrow x\in\{-1;0;1\}$. $S=\displaystyle\int_{-1}^{1}\left|x^{3} - x\right|\,\mathrm{d}x=\dfrac{1}{2}$.}]{Tính diện tích hình phẳng giới hạn bởi đồ thị các hàm số $y=x^{3} - x$ và $y=0$.}{0,5}
